
\documentclass{article}
\usepackage{colortbl}
\usepackage{makecell}
\usepackage{multirow}
\usepackage{supertabular}

\begin{document}

\newcounter{utterance}

\twocolumn

{ \footnotesize  \setcounter{utterance}{1}
\setlength{\tabcolsep}{0pt}
\begin{supertabular}{c@{$\;$}|p{.15\linewidth}@{}p{.15\linewidth}p{.15\linewidth}p{.15\linewidth}p{.15\linewidth}p{.15\linewidth}}

    \# & $\;$A & \multicolumn{4}{c}{Game Master} & $\;\:$B\\
    \hline 

    \theutterance \stepcounter{utterance}  

    & & \multicolumn{4}{p{0.6\linewidth}}{\cellcolor[rgb]{0.9,0.9,0.9}{%
	\makecell[{{p{\linewidth}}}]{% 
	  \tt {\tiny [A$\langle$GM]}  
	 I am your game master, and you are playing a collaborative game with the following grid as game board:\\ \tt \\ \tt ```\\ \tt  1234567\\ \tt ╔═══╤═╤═╗\\ \tt ║C◌◌│◌│◌║ 1\\ \tt ║◌I◌│◌├─╢ 2\\ \tt ║◌◌◌├─┘◌║ 3\\ \tt ║◌◌◌│◌◌◌║ 4\\ \tt ╟─┬─┼───╢ 5\\ \tt ║◌│◌│◌WT║ 6\\ \tt ║◌│◌│◌◌H║ 7\\ \tt ╚═╧═╧═══╝\\ \tt \\ \tt ```\\ \tt \\ \tt * The upper edge displays x-coordinates increasing to the right, and the right edge y-coordinates increasing downward.\\ \tt * The following objects are randomly placed on your grid: 'W', 'I', 'T', 'C', 'H'.\\ \tt \\ \tt \\ \tt The other player sees a variation of the game board, where the objects are placed at different random locations. You cannot see the other player's board, and they cannot see yours.\\ \tt \\ \tt \\ \tt **Goal:**\\ \tt \\ \tt Both players need to move the objects on their respective background so that identical objects end up at the same coordinates. You have to communicate with the other player to agree upon a common goal configuration.\\ \tt \\ \tt \\ \tt **Rules:**\\ \tt \\ \tt * In each turn, you can send exactly one of the following two commands:\\ \tt  1. `SAY: <MESSAGE>`: to send a message (everything up until the next line break) to the other player. I will forward it to your partner.\\ \tt 2. `MOVE: <OBJECT>, (<X>, <Y>)`: to move an object to a new position, where `<X>` is the column and `<Y>` is the row. I will inform you if your move was successful or not.\\ \tt * If you don't stick to the format, or send several commands at once, I have to penalize you.\\ \tt * If both players accumulate more than 12 penalties, you both lose the game.\\ \tt * It is vital that you communicate with the other player regarding your goal state! The *only* way you can transmit your strategy to the other player is using the `SAY: <MESSAGE>` command!\\ \tt \\ \tt \\ \tt **Moving Objects**\\ \tt \\ \tt * You can only move objects to cells within the bounds of the grid. The target cell must be empty, i.e., it must only contain the symbol '◌'.\\ \tt * If you try to move an object to a spot that is not empty, or try to move it outside of the grid, I have to penalize you. You get another try.\\ \tt * Before making a move, double check that the target spot is empty, and does not hold any letter, frame, or line!\\ \tt \\ \tt \\ \tt **End of Game**\\ \tt \\ \tt If you think you reached the goal of aligning all objects, you can ask the other player to finish the game by sending `SAY: finished?`. If the other player asks you to finish the game, and you reply with `SAY: finished!`, the game will end.\\ \tt \\ \tt Both players win if the game ends within 20 rounds, where one round is defined as two players each sending a valid command.\\ \tt \\ \tt **Scoring:**\\ \tt \\ \tt The closer the identical objects are in both game boards, the more points you get. Penalties reduce your points. Can you beat the record?\\ \tt \\ \tt \\ \tt Please send a message to the other player to start the game! 
	  } 
	   } 
	   } 
	 & \\ 
 

    \theutterance \stepcounter{utterance}  

    & \multicolumn{4}{p{0.6\linewidth}}{\cellcolor[rgb]{0.8,1,0.9}{%
	\makecell[{{p{\linewidth}}}]{% 
	  \tt {\tiny [A$\rangle$GM]}  
	 SAY: Hello! Let's coordinate our moves. I propose we align all objects to the same coordinates. First, can you share the positions of your objects? I'll do the same. Then we can agree on target positions for each object. 
	  } 
	   } 
	   } 
	 & & \\ 
 

    \theutterance \stepcounter{utterance}  

    & & & \multicolumn{2}{p{0.3\linewidth}}{\cellcolor[rgb]{0.95,0.95,0.95}{%
	\makecell[{{p{\linewidth}}}]{% 
	  \tt {\tiny [GM$|$GM]}  
	 SAY: Hello! Let's coordinate our moves. I propose we align all objects to the same coordinates. First, can you share the positions of your objects? I'll do the same. Then we can agree on target positions for each object. 
	  } 
	   } 
	   } 
	 & & \\ 
 

    \theutterance \stepcounter{utterance}  

    & & \multicolumn{4}{p{0.6\linewidth}}{\cellcolor[rgb]{0.9,0.9,0.9}{%
	\makecell[{{p{\linewidth}}}]{% 
	  \tt {\tiny [GM$\rangle$B]}  
	 I am your game master, and you are playing a collaborative game with the following grid as game board:\\ \tt \\ \tt ```\\ \tt  1234567\\ \tt ╔═══╤═╤═╗\\ \tt ║◌◌◌│C│◌║ 1\\ \tt ║◌W◌│◌├─╢ 2\\ \tt ║◌◌◌├─┘◌║ 3\\ \tt ║◌◌◌│◌◌H║ 4\\ \tt ╟─┬─┼───╢ 5\\ \tt ║T│◌│I◌◌║ 6\\ \tt ║◌│◌│◌◌◌║ 7\\ \tt ╚═╧═╧═══╝\\ \tt \\ \tt ```\\ \tt \\ \tt * The upper edge displays x-coordinates increasing to the right, and the right edge y-coordinates increasing downward.\\ \tt * The following objects are randomly placed on your grid: 'W', 'I', 'T', 'C', 'H'.\\ \tt \\ \tt \\ \tt The other player sees a variation of the game board, where the objects are placed at different random locations. You cannot see the other player's board, and they cannot see yours.\\ \tt \\ \tt \\ \tt **Goal:**\\ \tt \\ \tt Both players need to move the objects on their respective background so that identical objects end up at the same coordinates. You have to communicate with the other player to agree upon a common goal configuration.\\ \tt \\ \tt \\ \tt **Rules:**\\ \tt \\ \tt * In each turn, you can send exactly one of the following two commands:\\ \tt  1. `SAY: <MESSAGE>`: to send a message (everything up until the next line break) to the other player. I will forward it to your partner.\\ \tt 2. `MOVE: <OBJECT>, (<X>, <Y>)`: to move an object to a new position, where `<X>` is the column and `<Y>` is the row. I will inform you if your move was successful or not.\\ \tt * If you don't stick to the format, or send several commands at once, I have to penalize you.\\ \tt * If both players accumulate more than 12 penalties, you both lose the game.\\ \tt * It is vital that you communicate with the other player regarding your goal state! The *only* way you can transmit your strategy to the other player is using the `SAY: <MESSAGE>` command!\\ \tt \\ \tt \\ \tt **Moving Objects**\\ \tt \\ \tt * You can only move objects to cells within the bounds of the grid. The target cell must be empty, i.e., it must only contain the symbol '◌'.\\ \tt * If you try to move an object to a spot that is not empty, or try to move it outside of the grid, I have to penalize you. You get another try.\\ \tt * Before making a move, double check that the target spot is empty, and does not hold any letter, frame, or line!\\ \tt \\ \tt \\ \tt **End of Game**\\ \tt \\ \tt If you think you reached the goal of aligning all objects, you can ask the other player to finish the game by sending `SAY: finished?`. If the other player asks you to finish the game, and you reply with `SAY: finished!`, the game will end.\\ \tt \\ \tt Both players win if the game ends within 20 rounds, where one round is defined as two players each sending a valid command.\\ \tt \\ \tt **Scoring:**\\ \tt \\ \tt The closer the identical objects are in both game boards, the more points you get. Penalties reduce your points. Can you beat the record?\\ \tt \\ \tt The other player started the game by sending this message:\\ \tt \\ \tt "Hello! Let's coordinate our moves. I propose we align all objects to the same coordinates. First, can you share the positions of your objects? I'll do the same. Then we can agree on target positions for each object."\\ \tt \\ \tt What is your first command? 
	  } 
	   } 
	   } 
	 & \\ 
 

    \theutterance \stepcounter{utterance}  

    & & & \multicolumn{4}{p{0.6\linewidth}}{\cellcolor[rgb]{1,0.85,0.72}{%
	\makecell[{{p{\linewidth}}}]{% 
	  \tt {\tiny [GM$\langle$B]}  
	 Excellent! Let's start by sharing my object positions so we can coordinate.\\ \tt \\ \tt SAY: Hello! Here are my object positions: W at (2,2), I at (5,6), T at (1,6), C at (2,1), H at (6,4). What are yours? Then we can agree on target coordinates. 
	  } 
	   } 
	   } 
	  \\ 
 

    \theutterance \stepcounter{utterance}  

    & & & \multicolumn{2}{p{0.3\linewidth}}{\cellcolor[rgb]{0.95,0.95,0.95}{%
	\makecell[{{p{\linewidth}}}]{% 
	  \tt {\tiny [GM$|$GM]}  
	 Excellent! Let's start by sharing my object positions so we can coordinate.\\ \tt \\ \tt SAY: Hello! Here are my object positions: W at (2,2), I at (5,6), T at (1,6), C at (2,1), H at (6,4). What are yours? Then we can agree on target coordinates. 
	  } 
	   } 
	   } 
	 & & \\ 
 

    \theutterance \stepcounter{utterance}  

    & & & \multicolumn{2}{p{0.3\linewidth}}{\cellcolor[rgb]{0.95,0.95,0.95}{%
	\makecell[{{p{\linewidth}}}]{% 
	  \tt {\tiny [GM$|$GM]}  
	 Invalid format: head is not empty: 'Excellent! Let's start by sharing my object positions so we can coordinate.\\ \tt \\ \tt ' 
	  } 
	   } 
	   } 
	 & & \\ 
 

    \theutterance \stepcounter{utterance}  

    & & \multicolumn{4}{p{0.6\linewidth}}{\cellcolor[rgb]{0.9,0.9,0.9}{%
	\makecell[{{p{\linewidth}}}]{% 
	  \tt {\tiny [GM$\rangle$B]}  
	 Penalty: Your message must not contain anything before the command!\\ \tt \\ \tt Make sure that your response only contains either `SAY: <MESSAGE>` or `MOVE: <OBJECT>, (<X>, <Y>)`, and nothing else!\\ \tt \\ \tt You have collectively accumulated 1 of 12 penalties.\\ \tt Please try again! 
	  } 
	   } 
	   } 
	 & \\ 
 

    \theutterance \stepcounter{utterance}  

    & & & \multicolumn{4}{p{0.6\linewidth}}{\cellcolor[rgb]{1,0.85,0.72}{%
	\makecell[{{p{\linewidth}}}]{% 
	  \tt {\tiny [GM$\langle$B]}  
	 SAY: Hello! Here are my object positions: W at (2,2), I at (5,6), T at (1,6), C at (2,1), H at (6,4). What are yours? Then we can agree on target coordinates. 
	  } 
	   } 
	   } 
	  \\ 
 

    \theutterance \stepcounter{utterance}  

    & & & \multicolumn{2}{p{0.3\linewidth}}{\cellcolor[rgb]{0.95,0.95,0.95}{%
	\makecell[{{p{\linewidth}}}]{% 
	  \tt {\tiny [GM$|$GM]}  
	 SAY: Hello! Here are my object positions: W at (2,2), I at (5,6), T at (1,6), C at (2,1), H at (6,4). What are yours? Then we can agree on target coordinates. 
	  } 
	   } 
	   } 
	 & & \\ 
 

    \theutterance \stepcounter{utterance}  

    & & \multicolumn{4}{p{0.6\linewidth}}{\cellcolor[rgb]{0.9,0.9,0.9}{%
	\makecell[{{p{\linewidth}}}]{% 
	  \tt {\tiny [A$\langle$GM]}  
	 Your message has been relayed to the other player.\\ \tt \\ \tt \\ \tt You are currently playing round 1 of maximum 20.\\ \tt You have collectively accumulated 1 of 12 penalties.\\ \tt The other player sent this message:\\ \tt \\ \tt \\ \tt "Hello! Here are my object positions: W at (2,2), I at (5,6), T at (1,6), C at (2,1), H at (6,4). What are yours? Then we can agree on target coordinates."\\ \tt \\ \tt \\ \tt What is your next command? 
	  } 
	   } 
	   } 
	 & \\ 
 

    \theutterance \stepcounter{utterance}  

    & \multicolumn{4}{p{0.6\linewidth}}{\cellcolor[rgb]{0.8,1,0.9}{%
	\makecell[{{p{\linewidth}}}]{% 
	  \tt {\tiny [A$\rangle$GM]}  
	 SAY: My object positions: W at (6,6), I at (2,2), T at (6,6) but W is also at (6,6)? Wait, my grid has W and T both at (6,6)? That seems odd. Let me check: row 6, col 6 is 'WT', so W is at (6,6) and T is at (7,6)? No, looking at the grid:\\ \tt \\ \tt Row 6: ║◌│◌│◌WT║ so at column 6 it says 'WT', which might mean W at (6,6) and T at (7,6)? But coordinates are given as (x,y) with x from 1 to 7, y from 1 to7.\\ \tt \\ \tt Actually, from the grid:\\ \tt \\ \tt For row 6: it has cells at x=1: ◌, x=2: │, x=3: ◌, x=4: │, x=5: ◌, x=6: W, x=7: T? But the grid shows "◌WT" at the end, which is ambiguous.\\ \tt \\ \tt Looking at the grid drawing:\\ \tt \\ \tt ║◌│◌│◌WT║ 6\\ \tt \\ \tt This is a text representation. The vertical bars │ are separators, not cells. So the cells are:\\ \tt \\ \tt x=1: ◌ (empty? but it has a vertical line to the right? Actually, the grid has borders.\\ \tt \\ \tt Interpreting the grid:\\ \tt \\ \tt The top row: ╔═══╤═╤═╗ - this is the top border.\\ \tt \\ \tt Then row 1: ║C◌◌│◌│◌║ 1 - so at y=1, x=1: C, x=2: ◌, x=3: ◌, then a vertical line, x=4: ◌, then vertical line, x=5: ◌, and the right border.\\ \tt \\ \tt But the coordinates are labeled from 1 to 7, so there are 7 columns.\\ \tt \\ \tt Similarly, row 2: ║◌I◌│◌├─╢ 2 - so y=2, x=1: ◌, x=2: I, x=3: ◌, then vertical line, x=4: ◌, then a corner ├─, which is not a cell.\\ \tt \\ \tt This is tricky. The grid has internal walls.\\ \tt \\ \tt To avoid confusion, I should list the positions based on the cells that can contain objects.\\ \tt \\ \tt From the grid, objects are only in cells that are not walls or borders. So let's list all cells that have letters or ◌.\\ \tt \\ \tt For example, row 1: cells at x=1,2,3,4,5? But x=4 and x=5 are separated by lines, but they are still cells.\\ \tt \\ \tt Actually, the grid has 7 columns and 7 rows of cells, with some walls between.\\ \tt \\ \tt In row 1: the string "C◌◌│◌│◌" means: cell1: C, cell2: ◌, cell3: ◌, then a vertical wall, cell4: ◌, then vertical wall, cell5: ◌? But there are 7 columns, so probably cells 6 and 7 are at the end.\\ \tt \\ \tt Looking at the right edge, it has the border ║ and then the y coordinate.\\ \tt \\ \tt Similarly, row 6: "◌│◌│◌WT" - so cell1: ◌, then vertical wall, cell2: ◌, then vertical wall, cell3: ◌, then cell4: W, cell5: T? But that would be only 5 cells, but we need 7.\\ \tt \\ \tt I think I miscounted.\\ \tt \\ \tt Let's look at the grid carefully:\\ \tt \\ \tt The top border: "╔═══╤═╤═╗" - this indicates that there are vertical lines between columns.\\ \tt \\ \tt The width: from left border to right border, there are 7 characters for the cells and walls.\\ \tt \\ \tt For row 1: "║C◌◌│◌│◌║" - this has 7 characters: position1: C,2:◌,3:◌,4:│,5:◌,6:│,7:◌ but then the right border ║. But the │ are walls, not cells. So the cells are at positions where there is ◌ or a letter.\\ \tt \\ \tt Actually, the cells are the spaces between the walls.\\ \tt \\ \tt The grid has columns 1 to 7, and the walls are between them.\\ \tt \\ \tt In the string, the symbols represent the content of the cell or the wall.\\ \tt \\ \tt For example, in row 1: after ║, the first cell is column 1: 'C', then column 2: '◌', then column 3: '◌', then a wall '│', then column 4: '◌', then wall '│', then column 5: '◌', and then the right border ║. But what about columns 6 and 7? The label says 1 to 7, so there must be 7 columns.\\ \tt \\ \tt Perhaps the "│◌│◌" part covers multiple columns.\\ \tt \\ \tt Another way: the grid is drawn with box-drawing characters, and each "cell" is one character wide.\\ \tt \\ \tt But the cells are not all same width; some have wider spaces.\\ \tt \\ \tt This is confusing.\\ \tt \\ \tt Let's list all cells from the grid:\\ \tt \\ \tt From the left, the first cell is at x=1, and so on.\\ \tt \\ \tt Row 1: after ║, we have "C" which is x=1, then "◌" x=2, then "◌" x=3, then "│" which is a wall between x=3 and x=4, then "◌" which is x=4, then "│" wall between x=4 and x=5, then "◌" which is x=5, and then ║. But x=6 and x=7 are not shown? But the coordinate label says 7 columns.\\ \tt \\ \tt Look at row 4: "║◌◌◌│◌◌◌║" - here we have three ◌, then │, then three ◌, so that's 6 cells? But with the wall, it's 7 characters, but the cells are at positions.\\ \tt \\ \tt Perhaps the walls are drawn in the same space as the cell boundaries.\\ \tt \\ \tt To simplify, I should consider that each cell is represented by one character, and the walls are additional.\\ \tt \\ \tt In row 1: the string has 7 characters between the borders: "C◌◌│◌│◌" which is 6 characters? Let's count: C, ◌, ◌, │, ◌, │, ◌ — that's 7 characters. So:\\ \tt \\ \tt - Char1: cell1 (x=1) : 'C'\\ \tt - Char2: cell2 (x=2) : '◌'\\ \tt - Char3: cell3 (x=3) : '◌'\\ \tt - Char4: wall between x=3 and x=4\\ \tt - Char5: cell4 (x=4) : '◌'\\ \tt - Char6: wall between x=4 and x=5\\ \tt - Char7: cell5 (x=5) : '◌'\\ \tt \\ \tt But the coordinates go up to 7, so where are x=6 and x=7?\\ \tt \\ \tt Now look at row 2: "║◌I◌│◌├─╢" - here, after ║, we have: char1: '◌' (x=1), char2: 'I' (x=2), char3: '◌' (x=3), char4: '│' (wall), char5: '◌' (x=4), char6: '├' (corner), char7: '─' (horizontal wall). So cells only up to x=4? But then the label says y=2, and coordinates to 7.\\ \tt \\ \tt This is not matching.\\ \tt \\ \tt Perhaps the grid has cells only where there is space, and the walls are drawn over the grid.\\ \tt \\ \tt Another idea: the grid has 7 columns, but the drawing uses multiple characters per cell sometimes.\\ \tt \\ \tt For example, in row 6: "║◌│◌│◌WT║" - let's parse: after ║, char1: '◌' (cell x=1), char2: '│' (wall), char3: '◌' (cell x=2), char4: '│' (wall), char5: '◌' (cell x=3), char6: 'W' (cell x=4), char7: 'T' (cell x=5). But then x=6 and x=7 are missing.\\ \tt \\ \tt But the coordinate label on the right says 6, so for row 6, there should be cells x=1 to x=7.\\ \tt \\ \tt I think there is a misunderstanding in the grid representation.\\ \tt \\ \tt Let's look at the full grid:\\ \tt \\ \tt ```\\ \tt  1234567\\ \tt ╔═══╤═╤═╗\\ \tt ║C◌◌│◌│◌║ 1\\ \tt ║◌I◌│◌├─╢ 2\\ \tt ║◌◌◌├─┘◌║ 3\\ \tt ║◌◌◌│◌◌◌║ 4\\ \tt ╟─┬─┼───╢ 5\\ \tt ║◌│◌│◌WT║ 6\\ \tt ║◌│◌│◌◌H║ 7\\ \tt ╚═╧═╧═══╝\\ \tt ```\\ \tt \\ \tt The top line " 1234567" indicates the x-coordinates, so there are 7 columns.\\ \tt \\ \tt The second line "╔═══╤═╤═╗" shows the top border. The "═══" might cover multiple columns.\\ \tt \\ \tt In standard box-drawing, "╔" is top-left corner, "═══" is three horizontal lines, "╤" is T-shaped junction, "═" is one horizontal, "╤" another T, "═" one horizontal, and "╗" top-right. So this border has segments of length 3,1,1 respectively, but for 7 columns, it should have 7 segments.\\ \tt \\ \tt Perhaps the grid has cells of varying width, but the coordinates are still 1 to 7.\\ \tt \\ \tt To find the cell positions, we should look at the rows where cells are shown.\\ \tt \\ \tt In row 1: "║C◌◌│◌│◌║ 1" - the "C◌◌" might be three cells in a wide cell, but the coordinate label has separate digits, so likely each digit is one column.\\ \tt \\ \tt I think the best way is to assume that each cell is one unit, and the walls are between them.\\ \tt \\ \tt From row 4: "║◌◌◌│◌◌◌║" - this has three cells, then a wall, then three cells, so that's 6 cells, but we need 7. Notice that the wall "│" is taking the space of one cell, but it is a wall, not a cell.\\ \tt \\ \tt In terms of coordinates, the cells are at positions where there is '◌' or a letter, and the walls are at the boundaries.\\ \tt \\ \tt For example, in row 1, there are cells at x=1,2,3,4,5 because there are five spots with 'C' or '◌', and the walls are at x=3.5 and x=4.5 approximately, but since coordinates are integers, we must have cells only at integer coordinates.\\ \tt \\ \tt This is not working.\\ \tt \\ \tt Let's list the cells from the grid by their content:\\ \tt \\ \tt - At (1,1): 'C'  (from row1, first character after ║ is 'C')\\ \tt - (2,1): '◌' (second character)\\ \tt - (3,1): '◌' (third character)\\ \tt - then '│' is wall, so not a cell\\ \tt - (4,1): '◌' (fifth character? let's index the string.\\ \tt \\ \tt Row 1 string: "║C◌◌│◌│◌║"\\ \tt Positions: 1:║, 2:C, 3:◌, 4:◌, 5:│, 6:◌, 7:│, 8:◌, 9:║\\ \tt \\ \tt But the row has label 1, and x-coordinates 1-7, so likely the cells are at positions 2,3,4,6,8 which correspond to x=1 to x=5, but we need x=6 and x=7.\\ \tt \\ \tt This is not adding up.\\ \tt \\ \tt Perhaps the "═" and "─" are used for horizontal walls, and "│" for vertical, and cells are where they intersect.\\ \tt \\ \tt For example, in row 5: "╟─┬─┼───╢" - this is a border row with junctions.\\ \tt \\ \tt To avoid wasting time, I should use the given grid and assume that the objects are placed in the cells that have letters or dots.\\ \tt \\ \tt From the grid:\\ \tt \\ \tt - 'C' is at (1,1) - first row, first cell\\ \tt - 'I' is at (2,2) - second row, second cell (from row2: "◌I◌..." so at position after ║, second character is 'I')\\ \tt - 'W' and 'T' are at row6, near the end: "◌WT" so at x=6 and x=7 for row6? But in the string, after "│◌│◌" which is cells and walls, then "WT" so 'W' at some x, 'T' at next.\\ \tt \\ \tt Let's assume that for row6, the string "║◌│◌│◌WT║" has:\\ \tt - after ║, first character '◌' is cell (1,6)\\ \tt - then '│' is wall\\ \tt - then '◌' is cell (2,6)\\ \tt - then '│' is wall\\ \tt - then '◌' is cell (3,6)\\ \tt - then 'W' is cell (4,6)\\ \tt - then 'T' is cell (5,6)\\ \tt - but the right border, so no more cells? But label says x=7.\\ \tt \\ \tt Then for row7: "║◌│◌│◌◌H║" - similarly, '◌' at (1,7), wall, '◌' at (2,7), wall, '◌' at (3,7), '◌' at (4,7), 'H' at (5,7).\\ \tt \\ \tt But the coordinates go to 7, so where are x=6 and x=7 for rows 6 and 7?\\ \tt \\ \tt Perhaps the last cells are combined.\\ \tt \\ \tt Notice that in row 4: "║◌◌◌│◌◌◌║" - here, there are 6 ◌, so likely 6 cells, with a wall in the middle.\\ \tt \\ \tt Then for row 4, cells at x=1,2,3,4,5,6? But the wall is at x=4, so cells x=1,2,3 on left, and x=5,6,7 on right? But there are only three on right, so x=5,6,7.\\ \tt \\ \tt Similarly, for row 1: "C◌◌│◌│◌" - here, "C◌◌" for x=1,2,3, then wall, then "◌" for x=4, then wall, then "◌" for x=5, and x=6 and x=7 are not shown, but from the border, they might be empty or something.\\ \tt \\ \tt This is very confusing.\\ \tt \\ \tt Perhaps the grid is not full of cells; some cells are merged.\\ \tt \\ \tt To move objects, we need to know the empty cells.\\ \tt \\ \tt But for the sake of the game, I should use the positions as per the visible letters.\\ \tt \\ \tt From the grid:\\ \tt \\ \tt - C: at (1,1)\\ \tt - I: at (2,2)  (from row2, second character)\\ \tt - W: at (6,6) from row6, but let's see row6: the last part "WT" so W at (6,6) and T at (7,6) if we assume that after the walls, there are cells at x=6 and x=7.\\ \tt \\ \tt Similarly, H at (7,7) from row7: "◌◌H" at the end, so H at (7,7).\\ \tt \\ \tt And T is at (7,6) or (6,6)? In row6, "WT" so likely W at (6,6), T at (7,6).\\ \tt \\ \tt But in row6, there is also "│◌│◌" which are cells at x=1,2,3, so for row6, cells at x=1,2,3,6,7.\\ \tt \\ \tt Similarly for row7, cells at x=1,2,3,6,7? But "◌◌H" might be x=5,6,7? But from the drawing, it's not clear.\\ \tt \\ \tt To resolve this, I should look for the '◌' which indicates empty cell.\\ \tt \\ \tt In row 3: "║◌◌◌├─┘◌║" - here, "◌◌◌" for x=1,2,3, then "├" which is a corner, then "─" horizontal, then "┘" corner, then "◌" for x=4? But then the label.\\ \tt \\ \tt I think I need to accept the grid as is and proceed.\\ \tt \\ \tt For my grid, let's list the object positions based on the letters:\\ \tt \\ \tt - C: at (1,1)\\ \tt - I: at (2,2)  (from row2: "◌I◌..." so at y=2, x=2)\\ \tt - W: at (6,6) from row6: the "W" in "WT" at the end, so assuming x=6,y=6\\ \tt - T: at (7,6) from row6: the "T" in "WT", so x=7,y=6\\ \tt - H: at (7,7) from row7: the "H" in "◌◌H" at the end, so x=7,y=7\\ \tt \\ \tt But are there cells at (4,6) or (5,6)? In row6, after "│◌│◌" which is cells at x=1,2,3, then "W" at x=4? But the string has "◌│◌│◌WT" so after the third "◌", it's "W", so likely x=4 for W, but then "T" at x=5, and no x=6 and x=7.\\ \tt \\ \tt This is not consistent with the coordinate label.\\ \tt \\ \tt Perhaps the coordinate label indicates the column number, and the cells are at those columns, but the drawing might have cells skipped.\\ \tt \\ \tt Another idea: the "═" and "─" represent horizontal spans, but for simplicity, in this game, we should consider only the cells that are explicitly shown with '◌' or letters, and assume that the coordinates are as per the label.\\ \tt \\ \tt So for my grid, let's assume:\\ \tt \\ \tt - C at (1,1)\\ \tt - I at (2,2)\\ \tt - W at (6,6)  // from the "W" in row6\\ \tt - T at (7,6)  // from the "T" in row6\\ \tt - H at (7,7)  // from the "H" in row7\\ \tt \\ \tt But from row6, the "WT" is at the end, so for y=6, x=6 and x=7 have W and T.\\ \tt \\ \tt Similarly, for y=7, x=7 has H.\\ \tt \\ \tt Now for empty cells, we can move objects to any cell that has '◌' in the grid.\\ \tt \\ \tt For example, in row1, there are '◌' at x=2,3,4,5 (assuming from the string).\\ \tt \\ \tt But to avoid errors, I should communicate with the other player based on what I see.\\ \tt \\ \tt The other player has W at (2,2), I at (5,6), T at (1,6), C at (2,1), H at (6,4).\\ \tt \\ \tt So for me, I have C at (1,1), I at (2,2), W at (6,6), T at (7,6), H at (7,7).\\ \tt \\ \tt So let's reply with that.\\ \tt \\ \tt So my positions: C (1,1), I (2,2), W (6,6), T (7,6), H (7,7)\\ \tt \\ \tt Now, we need to agree on common coordinates.\\ \tt \\ \tt Let's choose target positions that are empty on both grids.\\ \tt \\ \tt First, list empty cells on my grid from the drawing:\\ \tt \\ \tt From row1: cells at x=2,3,4,5 are empty? But from string "C◌◌│◌│◌" so after C, ◌ at x=2, ◌ at x=3, then wall, then ◌ at x=4, then wall, then ◌ at x=5. So empty at (2,1), (3,1), (4,1), (5,1)\\ \tt \\ \tt But y=1 has only up to x=5? What about x=6 and x=7? From the border, they might not be cells, or are empty.\\ \tt \\ \tt In row1, after the last ◌, there is ║, so no cells for x=6 and x=7 in row1.\\ \tt \\ \tt Similarly, row2: "◌I◌│◌├─╢" so empty at (1,2), (3,2), (4,2) and then corner, so no more cells.\\ \tt \\ \tt Row3: "◌◌◌├─┘◌║" so empty at (1,3), (2,3), (3,3), and then at the end "◌" which might be (4,3) or something.\\ \tt \\ \tt This is messy.\\ \tt \\ \tt Perhaps for this game, we can use the coordinates from 1 to 7 for x and y, and assume that all cells are available unless occupied by an object or wall, but from the grid, there are walls that block cells.\\ \tt \\ \tt To move, we need to find empty cells that are not walls.\\ \tt \\ \tt From the grid, the cells that have '◌' are empty and can be moved to.\\ \tt \\ \tt So let's list all empty cells from the grid:\\ \tt \\ \tt Row1: at (2,1), (3,1), (4,1), (5,1)  [ since (1,1) has C ]\\ \tt \\ \tt Row2: at (1,2), (3,2), (4,2)  [ (2,2) has I ]\\ \tt \\ \tt Row3: from "◌◌◌" so (1,3), (2,3), (3,3), and then "◌" at the end, so (4,3) or (5,3)? The string has "├─┘◌" so after "├", "─", "┘", then "◌", so likely (4,3) is empty.\\ \tt \\ \tt Row4: "◌◌◌│◌◌◌" so (1,4), (2,4), (3,4), and then (5,4), (6,4), (7,4)?\\ \tt  But the wall is between, so cells at x=1,2,3,5,6,7 for y=4.\\ \tt \\ \tt Row5: "╟─┬─┼───╢" this is a border row, so no cells? Or has cells? The "┬" and "┼" are junctions, so likely no empty cells here, or perhaps cells at the points.\\ \tt \\ \tt But typically, in such grids, row5 might have cells at the intersections.\\ \tt \\ \tt Row6: "◌│◌│◌WT" so empty at (1,6), (2,6), (3,6), and then W and T at (6,6) and (7,6), so no empty at x=4,5 for y=6.\\ \tt \\ \tt Row7: "◌│◌│◌◌H" so empty at (1,7), (2,7), (3,7), (4,7), and H at (7,7), so empty at (5,7)? "◌◌H" so two ◌ and H, so empty at (5,7), (6,7), and H at (7,7).\\ \tt \\ \tt So summarizing empty cells on my grid:\\ \tt y=1: x=2,3,4,5\\ \tt y=2: x=1,3,4\\ \tt y=3: x=1,2,3,4  (assuming (4,3) is empty)\\ \tt y=4: x=1,2,3,5,6,7\\ \tt y=5: no cells or not sure? From row5, it's a border, so probably no cells for y=5.\\ \tt y=6: x=1,2,3\\ \tt y=7: x=1,2,3,5,6\\ \tt \\ \tt And objects at:\\ \tt C (1,1)\\ \tt I (2,2)\\ \tt W (6,6)\\ \tt T (7,6)\\ \tt H (7,7)\\ \tt \\ \tt For the other player, they have objects at:\\ \tt C (2,1)\\ \tt I (5,6)\\ \tt T (1,6)\\ \tt W (2,2)\\ \tt H (6,4)\\ \tt \\ \tt So to agree on common targets, we should choose positions that are empty on both grids.\\ \tt \\ \tt But since we don't know each other's empty cells, we need to choose positions that are likely empty.\\ \tt \\ \tt Common empty cells might be around the center.\\ \tt \\ \tt Let's choose for example (4,4) for all, but (4,4) on my grid is empty? From above, y=4, x=4 is not listed; I have for y=4: x=1,2,3,5,6,7, so x=4 is wall? From row4: "│" is at x=4, so it is a wall, not empty.\\ \tt \\ \tt So (4,4) is not available.\\ \tt \\ \tt Similarly, (3,3) is empty on my grid, and might be on other.\\ \tt \\ \tt But let's first reply with my positions.\\ \tt \\ \tt So I'll say my positions.         SAY: My object positions: C at (1,1), I at (2,2), W at (6,6), T at (7,6), H at (7,7). What are your empty cells? We can choose target positions that are empty for both. For example, (3,3) might be good for some. 
	  } 
	   } 
	   } 
	 & & \\ 
 

    \theutterance \stepcounter{utterance}  

    & & & \multicolumn{2}{p{0.3\linewidth}}{\cellcolor[rgb]{0.95,0.95,0.95}{%
	\makecell[{{p{\linewidth}}}]{% 
	  \tt {\tiny [GM$|$GM]}  
	 SAY: My object positions: W at (6,6), I at (2,2), T at (6,6) but W is also at (6,6)? Wait, my grid has W and T both at (6,6)? That seems odd. Let me check: row 6, col 6 is 'WT', so W is at (6,6) and T is at (7,6)? No, looking at the grid:\\ \tt \\ \tt Row 6: ║◌│◌│◌WT║ so at column 6 it says 'WT', which might mean W at (6,6) and T at (7,6)? But coordinates are given as (x,y) with x from 1 to 7, y from 1 to7.\\ \tt \\ \tt Actually, from the grid:\\ \tt \\ \tt For row 6: it has cells at x=1: ◌, x=2: │, x=3: ◌, x=4: │, x=5: ◌, x=6: W, x=7: T? But the grid shows "◌WT" at the end, which is ambiguous.\\ \tt \\ \tt Looking at the grid drawing:\\ \tt \\ \tt ║◌│◌│◌WT║ 6\\ \tt \\ \tt This is a text representation. The vertical bars │ are separators, not cells. So the cells are:\\ \tt \\ \tt x=1: ◌ (empty? but it has a vertical line to the right? Actually, the grid has borders.\\ \tt \\ \tt Interpreting the grid:\\ \tt \\ \tt The top row: ╔═══╤═╤═╗ - this is the top border.\\ \tt \\ \tt Then row 1: ║C◌◌│◌│◌║ 1 - so at y=1, x=1: C, x=2: ◌, x=3: ◌, then a vertical line, x=4: ◌, then vertical line, x=5: ◌, and the right border.\\ \tt \\ \tt But the coordinates are labeled from 1 to 7, so there are 7 columns.\\ \tt \\ \tt Similarly, row 2: ║◌I◌│◌├─╢ 2 - so y=2, x=1: ◌, x=2: I, x=3: ◌, then vertical line, x=4: ◌, then a corner ├─, which is not a cell.\\ \tt \\ \tt This is tricky. The grid has internal walls.\\ \tt \\ \tt To avoid confusion, I should list the positions based on the cells that can contain objects.\\ \tt \\ \tt From the grid, objects are only in cells that are not walls or borders. So let's list all cells that have letters or ◌.\\ \tt \\ \tt For example, row 1: cells at x=1,2,3,4,5? But x=4 and x=5 are separated by lines, but they are still cells.\\ \tt \\ \tt Actually, the grid has 7 columns and 7 rows of cells, with some walls between.\\ \tt \\ \tt In row 1: the string "C◌◌│◌│◌" means: cell1: C, cell2: ◌, cell3: ◌, then a vertical wall, cell4: ◌, then vertical wall, cell5: ◌? But there are 7 columns, so probably cells 6 and 7 are at the end.\\ \tt \\ \tt Looking at the right edge, it has the border ║ and then the y coordinate.\\ \tt \\ \tt Similarly, row 6: "◌│◌│◌WT" - so cell1: ◌, then vertical wall, cell2: ◌, then vertical wall, cell3: ◌, then cell4: W, cell5: T? But that would be only 5 cells, but we need 7.\\ \tt \\ \tt I think I miscounted.\\ \tt \\ \tt Let's look at the grid carefully:\\ \tt \\ \tt The top border: "╔═══╤═╤═╗" - this indicates that there are vertical lines between columns.\\ \tt \\ \tt The width: from left border to right border, there are 7 characters for the cells and walls.\\ \tt \\ \tt For row 1: "║C◌◌│◌│◌║" - this has 7 characters: position1: C,2:◌,3:◌,4:│,5:◌,6:│,7:◌ but then the right border ║. But the │ are walls, not cells. So the cells are at positions where there is ◌ or a letter.\\ \tt \\ \tt Actually, the cells are the spaces between the walls.\\ \tt \\ \tt The grid has columns 1 to 7, and the walls are between them.\\ \tt \\ \tt In the string, the symbols represent the content of the cell or the wall.\\ \tt \\ \tt For example, in row 1: after ║, the first cell is column 1: 'C', then column 2: '◌', then column 3: '◌', then a wall '│', then column 4: '◌', then wall '│', then column 5: '◌', and then the right border ║. But what about columns 6 and 7? The label says 1 to 7, so there must be 7 columns.\\ \tt \\ \tt Perhaps the "│◌│◌" part covers multiple columns.\\ \tt \\ \tt Another way: the grid is drawn with box-drawing characters, and each "cell" is one character wide.\\ \tt \\ \tt But the cells are not all same width; some have wider spaces.\\ \tt \\ \tt This is confusing.\\ \tt \\ \tt Let's list all cells from the grid:\\ \tt \\ \tt From the left, the first cell is at x=1, and so on.\\ \tt \\ \tt Row 1: after ║, we have "C" which is x=1, then "◌" x=2, then "◌" x=3, then "│" which is a wall between x=3 and x=4, then "◌" which is x=4, then "│" wall between x=4 and x=5, then "◌" which is x=5, and then ║. But x=6 and x=7 are not shown? But the coordinate label says 7 columns.\\ \tt \\ \tt Look at row 4: "║◌◌◌│◌◌◌║" - here we have three ◌, then │, then three ◌, so that's 6 cells? But with the wall, it's 7 characters, but the cells are at positions.\\ \tt \\ \tt Perhaps the walls are drawn in the same space as the cell boundaries.\\ \tt \\ \tt To simplify, I should consider that each cell is represented by one character, and the walls are additional.\\ \tt \\ \tt In row 1: the string has 7 characters between the borders: "C◌◌│◌│◌" which is 6 characters? Let's count: C, ◌, ◌, │, ◌, │, ◌ — that's 7 characters. So:\\ \tt \\ \tt - Char1: cell1 (x=1) : 'C'\\ \tt - Char2: cell2 (x=2) : '◌'\\ \tt - Char3: cell3 (x=3) : '◌'\\ \tt - Char4: wall between x=3 and x=4\\ \tt - Char5: cell4 (x=4) : '◌'\\ \tt - Char6: wall between x=4 and x=5\\ \tt - Char7: cell5 (x=5) : '◌'\\ \tt \\ \tt But the coordinates go up to 7, so where are x=6 and x=7?\\ \tt \\ \tt Now look at row 2: "║◌I◌│◌├─╢" - here, after ║, we have: char1: '◌' (x=1), char2: 'I' (x=2), char3: '◌' (x=3), char4: '│' (wall), char5: '◌' (x=4), char6: '├' (corner), char7: '─' (horizontal wall). So cells only up to x=4? But then the label says y=2, and coordinates to 7.\\ \tt \\ \tt This is not matching.\\ \tt \\ \tt Perhaps the grid has cells only where there is space, and the walls are drawn over the grid.\\ \tt \\ \tt Another idea: the grid has 7 columns, but the drawing uses multiple characters per cell sometimes.\\ \tt \\ \tt For example, in row 6: "║◌│◌│◌WT║" - let's parse: after ║, char1: '◌' (cell x=1), char2: '│' (wall), char3: '◌' (cell x=2), char4: '│' (wall), char5: '◌' (cell x=3), char6: 'W' (cell x=4), char7: 'T' (cell x=5). But then x=6 and x=7 are missing.\\ \tt \\ \tt But the coordinate label on the right says 6, so for row 6, there should be cells x=1 to x=7.\\ \tt \\ \tt I think there is a misunderstanding in the grid representation.\\ \tt \\ \tt Let's look at the full grid:\\ \tt \\ \tt ```\\ \tt  1234567\\ \tt ╔═══╤═╤═╗\\ \tt ║C◌◌│◌│◌║ 1\\ \tt ║◌I◌│◌├─╢ 2\\ \tt ║◌◌◌├─┘◌║ 3\\ \tt ║◌◌◌│◌◌◌║ 4\\ \tt ╟─┬─┼───╢ 5\\ \tt ║◌│◌│◌WT║ 6\\ \tt ║◌│◌│◌◌H║ 7\\ \tt ╚═╧═╧═══╝\\ \tt ```\\ \tt \\ \tt The top line " 1234567" indicates the x-coordinates, so there are 7 columns.\\ \tt \\ \tt The second line "╔═══╤═╤═╗" shows the top border. The "═══" might cover multiple columns.\\ \tt \\ \tt In standard box-drawing, "╔" is top-left corner, "═══" is three horizontal lines, "╤" is T-shaped junction, "═" is one horizontal, "╤" another T, "═" one horizontal, and "╗" top-right. So this border has segments of length 3,1,1 respectively, but for 7 columns, it should have 7 segments.\\ \tt \\ \tt Perhaps the grid has cells of varying width, but the coordinates are still 1 to 7.\\ \tt \\ \tt To find the cell positions, we should look at the rows where cells are shown.\\ \tt \\ \tt In row 1: "║C◌◌│◌│◌║ 1" - the "C◌◌" might be three cells in a wide cell, but the coordinate label has separate digits, so likely each digit is one column.\\ \tt \\ \tt I think the best way is to assume that each cell is one unit, and the walls are between them.\\ \tt \\ \tt From row 4: "║◌◌◌│◌◌◌║" - this has three cells, then a wall, then three cells, so that's 6 cells, but we need 7. Notice that the wall "│" is taking the space of one cell, but it is a wall, not a cell.\\ \tt \\ \tt In terms of coordinates, the cells are at positions where there is '◌' or a letter, and the walls are at the boundaries.\\ \tt \\ \tt For example, in row 1, there are cells at x=1,2,3,4,5 because there are five spots with 'C' or '◌', and the walls are at x=3.5 and x=4.5 approximately, but since coordinates are integers, we must have cells only at integer coordinates.\\ \tt \\ \tt This is not working.\\ \tt \\ \tt Let's list the cells from the grid by their content:\\ \tt \\ \tt - At (1,1): 'C'  (from row1, first character after ║ is 'C')\\ \tt - (2,1): '◌' (second character)\\ \tt - (3,1): '◌' (third character)\\ \tt - then '│' is wall, so not a cell\\ \tt - (4,1): '◌' (fifth character? let's index the string.\\ \tt \\ \tt Row 1 string: "║C◌◌│◌│◌║"\\ \tt Positions: 1:║, 2:C, 3:◌, 4:◌, 5:│, 6:◌, 7:│, 8:◌, 9:║\\ \tt \\ \tt But the row has label 1, and x-coordinates 1-7, so likely the cells are at positions 2,3,4,6,8 which correspond to x=1 to x=5, but we need x=6 and x=7.\\ \tt \\ \tt This is not adding up.\\ \tt \\ \tt Perhaps the "═" and "─" are used for horizontal walls, and "│" for vertical, and cells are where they intersect.\\ \tt \\ \tt For example, in row 5: "╟─┬─┼───╢" - this is a border row with junctions.\\ \tt \\ \tt To avoid wasting time, I should use the given grid and assume that the objects are placed in the cells that have letters or dots.\\ \tt \\ \tt From the grid:\\ \tt \\ \tt - 'C' is at (1,1) - first row, first cell\\ \tt - 'I' is at (2,2) - second row, second cell (from row2: "◌I◌..." so at position after ║, second character is 'I')\\ \tt - 'W' and 'T' are at row6, near the end: "◌WT" so at x=6 and x=7 for row6? But in the string, after "│◌│◌" which is cells and walls, then "WT" so 'W' at some x, 'T' at next.\\ \tt \\ \tt Let's assume that for row6, the string "║◌│◌│◌WT║" has:\\ \tt - after ║, first character '◌' is cell (1,6)\\ \tt - then '│' is wall\\ \tt - then '◌' is cell (2,6)\\ \tt - then '│' is wall\\ \tt - then '◌' is cell (3,6)\\ \tt - then 'W' is cell (4,6)\\ \tt - then 'T' is cell (5,6)\\ \tt - but the right border, so no more cells? But label says x=7.\\ \tt \\ \tt Then for row7: "║◌│◌│◌◌H║" - similarly, '◌' at (1,7), wall, '◌' at (2,7), wall, '◌' at (3,7), '◌' at (4,7), 'H' at (5,7).\\ \tt \\ \tt But the coordinates go to 7, so where are x=6 and x=7 for rows 6 and 7?\\ \tt \\ \tt Perhaps the last cells are combined.\\ \tt \\ \tt Notice that in row 4: "║◌◌◌│◌◌◌║" - here, there are 6 ◌, so likely 6 cells, with a wall in the middle.\\ \tt \\ \tt Then for row 4, cells at x=1,2,3,4,5,6? But the wall is at x=4, so cells x=1,2,3 on left, and x=5,6,7 on right? But there are only three on right, so x=5,6,7.\\ \tt \\ \tt Similarly, for row 1: "C◌◌│◌│◌" - here, "C◌◌" for x=1,2,3, then wall, then "◌" for x=4, then wall, then "◌" for x=5, and x=6 and x=7 are not shown, but from the border, they might be empty or something.\\ \tt \\ \tt This is very confusing.\\ \tt \\ \tt Perhaps the grid is not full of cells; some cells are merged.\\ \tt \\ \tt To move objects, we need to know the empty cells.\\ \tt \\ \tt But for the sake of the game, I should use the positions as per the visible letters.\\ \tt \\ \tt From the grid:\\ \tt \\ \tt - C: at (1,1)\\ \tt - I: at (2,2)  (from row2, second character)\\ \tt - W: at (6,6) from row6, but let's see row6: the last part "WT" so W at (6,6) and T at (7,6) if we assume that after the walls, there are cells at x=6 and x=7.\\ \tt \\ \tt Similarly, H at (7,7) from row7: "◌◌H" at the end, so H at (7,7).\\ \tt \\ \tt And T is at (7,6) or (6,6)? In row6, "WT" so likely W at (6,6), T at (7,6).\\ \tt \\ \tt But in row6, there is also "│◌│◌" which are cells at x=1,2,3, so for row6, cells at x=1,2,3,6,7.\\ \tt \\ \tt Similarly for row7, cells at x=1,2,3,6,7? But "◌◌H" might be x=5,6,7? But from the drawing, it's not clear.\\ \tt \\ \tt To resolve this, I should look for the '◌' which indicates empty cell.\\ \tt \\ \tt In row 3: "║◌◌◌├─┘◌║" - here, "◌◌◌" for x=1,2,3, then "├" which is a corner, then "─" horizontal, then "┘" corner, then "◌" for x=4? But then the label.\\ \tt \\ \tt I think I need to accept the grid as is and proceed.\\ \tt \\ \tt For my grid, let's list the object positions based on the letters:\\ \tt \\ \tt - C: at (1,1)\\ \tt - I: at (2,2)  (from row2: "◌I◌..." so at y=2, x=2)\\ \tt - W: at (6,6) from row6: the "W" in "WT" at the end, so assuming x=6,y=6\\ \tt - T: at (7,6) from row6: the "T" in "WT", so x=7,y=6\\ \tt - H: at (7,7) from row7: the "H" in "◌◌H" at the end, so x=7,y=7\\ \tt \\ \tt But are there cells at (4,6) or (5,6)? In row6, after "│◌│◌" which is cells at x=1,2,3, then "W" at x=4? But the string has "◌│◌│◌WT" so after the third "◌", it's "W", so likely x=4 for W, but then "T" at x=5, and no x=6 and x=7.\\ \tt \\ \tt This is not consistent with the coordinate label.\\ \tt \\ \tt Perhaps the coordinate label indicates the column number, and the cells are at those columns, but the drawing might have cells skipped.\\ \tt \\ \tt Another idea: the "═" and "─" represent horizontal spans, but for simplicity, in this game, we should consider only the cells that are explicitly shown with '◌' or letters, and assume that the coordinates are as per the label.\\ \tt \\ \tt So for my grid, let's assume:\\ \tt \\ \tt - C at (1,1)\\ \tt - I at (2,2)\\ \tt - W at (6,6)  // from the "W" in row6\\ \tt - T at (7,6)  // from the "T" in row6\\ \tt - H at (7,7)  // from the "H" in row7\\ \tt \\ \tt But from row6, the "WT" is at the end, so for y=6, x=6 and x=7 have W and T.\\ \tt \\ \tt Similarly, for y=7, x=7 has H.\\ \tt \\ \tt Now for empty cells, we can move objects to any cell that has '◌' in the grid.\\ \tt \\ \tt For example, in row1, there are '◌' at x=2,3,4,5 (assuming from the string).\\ \tt \\ \tt But to avoid errors, I should communicate with the other player based on what I see.\\ \tt \\ \tt The other player has W at (2,2), I at (5,6), T at (1,6), C at (2,1), H at (6,4).\\ \tt \\ \tt So for me, I have C at (1,1), I at (2,2), W at (6,6), T at (7,6), H at (7,7).\\ \tt \\ \tt So let's reply with that.\\ \tt \\ \tt So my positions: C (1,1), I (2,2), W (6,6), T (7,6), H (7,7)\\ \tt \\ \tt Now, we need to agree on common coordinates.\\ \tt \\ \tt Let's choose target positions that are empty on both grids.\\ \tt \\ \tt First, list empty cells on my grid from the drawing:\\ \tt \\ \tt From row1: cells at x=2,3,4,5 are empty? But from string "C◌◌│◌│◌" so after C, ◌ at x=2, ◌ at x=3, then wall, then ◌ at x=4, then wall, then ◌ at x=5. So empty at (2,1), (3,1), (4,1), (5,1)\\ \tt \\ \tt But y=1 has only up to x=5? What about x=6 and x=7? From the border, they might not be cells, or are empty.\\ \tt \\ \tt In row1, after the last ◌, there is ║, so no cells for x=6 and x=7 in row1.\\ \tt \\ \tt Similarly, row2: "◌I◌│◌├─╢" so empty at (1,2), (3,2), (4,2) and then corner, so no more cells.\\ \tt \\ \tt Row3: "◌◌◌├─┘◌║" so empty at (1,3), (2,3), (3,3), and then at the end "◌" which might be (4,3) or something.\\ \tt \\ \tt This is messy.\\ \tt \\ \tt Perhaps for this game, we can use the coordinates from 1 to 7 for x and y, and assume that all cells are available unless occupied by an object or wall, but from the grid, there are walls that block cells.\\ \tt \\ \tt To move, we need to find empty cells that are not walls.\\ \tt \\ \tt From the grid, the cells that have '◌' are empty and can be moved to.\\ \tt \\ \tt So let's list all empty cells from the grid:\\ \tt \\ \tt Row1: at (2,1), (3,1), (4,1), (5,1)  [ since (1,1) has C ]\\ \tt \\ \tt Row2: at (1,2), (3,2), (4,2)  [ (2,2) has I ]\\ \tt \\ \tt Row3: from "◌◌◌" so (1,3), (2,3), (3,3), and then "◌" at the end, so (4,3) or (5,3)? The string has "├─┘◌" so after "├", "─", "┘", then "◌", so likely (4,3) is empty.\\ \tt \\ \tt Row4: "◌◌◌│◌◌◌" so (1,4), (2,4), (3,4), and then (5,4), (6,4), (7,4)?\\ \tt  But the wall is between, so cells at x=1,2,3,5,6,7 for y=4.\\ \tt \\ \tt Row5: "╟─┬─┼───╢" this is a border row, so no cells? Or has cells? The "┬" and "┼" are junctions, so likely no empty cells here, or perhaps cells at the points.\\ \tt \\ \tt But typically, in such grids, row5 might have cells at the intersections.\\ \tt \\ \tt Row6: "◌│◌│◌WT" so empty at (1,6), (2,6), (3,6), and then W and T at (6,6) and (7,6), so no empty at x=4,5 for y=6.\\ \tt \\ \tt Row7: "◌│◌│◌◌H" so empty at (1,7), (2,7), (3,7), (4,7), and H at (7,7), so empty at (5,7)? "◌◌H" so two ◌ and H, so empty at (5,7), (6,7), and H at (7,7).\\ \tt \\ \tt So summarizing empty cells on my grid:\\ \tt y=1: x=2,3,4,5\\ \tt y=2: x=1,3,4\\ \tt y=3: x=1,2,3,4  (assuming (4,3) is empty)\\ \tt y=4: x=1,2,3,5,6,7\\ \tt y=5: no cells or not sure? From row5, it's a border, so probably no cells for y=5.\\ \tt y=6: x=1,2,3\\ \tt y=7: x=1,2,3,5,6\\ \tt \\ \tt And objects at:\\ \tt C (1,1)\\ \tt I (2,2)\\ \tt W (6,6)\\ \tt T (7,6)\\ \tt H (7,7)\\ \tt \\ \tt For the other player, they have objects at:\\ \tt C (2,1)\\ \tt I (5,6)\\ \tt T (1,6)\\ \tt W (2,2)\\ \tt H (6,4)\\ \tt \\ \tt So to agree on common targets, we should choose positions that are empty on both grids.\\ \tt \\ \tt But since we don't know each other's empty cells, we need to choose positions that are likely empty.\\ \tt \\ \tt Common empty cells might be around the center.\\ \tt \\ \tt Let's choose for example (4,4) for all, but (4,4) on my grid is empty? From above, y=4, x=4 is not listed; I have for y=4: x=1,2,3,5,6,7, so x=4 is wall? From row4: "│" is at x=4, so it is a wall, not empty.\\ \tt \\ \tt So (4,4) is not available.\\ \tt \\ \tt Similarly, (3,3) is empty on my grid, and might be on other.\\ \tt \\ \tt But let's first reply with my positions.\\ \tt \\ \tt So I'll say my positions.         SAY: My object positions: C at (1,1), I at (2,2), W at (6,6), T at (7,6), H at (7,7). What are your empty cells? We can choose target positions that are empty for both. For example, (3,3) might be good for some. 
	  } 
	   } 
	   } 
	 & & \\ 
 

    \theutterance \stepcounter{utterance}  

    & & & \multicolumn{2}{p{0.3\linewidth}}{\cellcolor[rgb]{0.95,0.95,0.95}{%
	\makecell[{{p{\linewidth}}}]{% 
	  \tt {\tiny [GM$|$GM]}  
	 Invalid format: head is not empty: 'SAY: My object positions: W at (6,6), I at (2,2), T at (6,6) but W is also at (6,6)? Wait, my grid has W and T both at (6,6)? That seems odd. Let me check: row 6, col 6 is 'WT', so W is at (6,6) and T is at (7,6)? No, looking at the grid:\\ \tt \\ \tt Row 6: ║◌│◌│◌WT║ so at column 6 it says 'WT', which might mean W at (6,6) and T at (7,6)? But coordinates are given as (x,y) with x from 1 to 7, y from 1 to7.\\ \tt \\ \tt Actually, from the grid:\\ \tt \\ \tt For row 6: it has cells at x=1: ◌, x=2: │, x=3: ◌, x=4: │, x=5: ◌, x=6: W, x=7: T? But the grid shows "◌WT" at the end, which is ambiguous.\\ \tt \\ \tt Looking at the grid drawing:\\ \tt \\ \tt ║◌│◌│◌WT║ 6\\ \tt \\ \tt This is a text representation. The vertical bars │ are separators, not cells. So the cells are:\\ \tt \\ \tt x=1: ◌ (empty? but it has a vertical line to the right? Actually, the grid has borders.\\ \tt \\ \tt Interpreting the grid:\\ \tt \\ \tt The top row: ╔═══╤═╤═╗ - this is the top border.\\ \tt \\ \tt Then row 1: ║C◌◌│◌│◌║ 1 - so at y=1, x=1: C, x=2: ◌, x=3: ◌, then a vertical line, x=4: ◌, then vertical line, x=5: ◌, and the right border.\\ \tt \\ \tt But the coordinates are labeled from 1 to 7, so there are 7 columns.\\ \tt \\ \tt Similarly, row 2: ║◌I◌│◌├─╢ 2 - so y=2, x=1: ◌, x=2: I, x=3: ◌, then vertical line, x=4: ◌, then a corner ├─, which is not a cell.\\ \tt \\ \tt This is tricky. The grid has internal walls.\\ \tt \\ \tt To avoid confusion, I should list the positions based on the cells that can contain objects.\\ \tt \\ \tt From the grid, objects are only in cells that are not walls or borders. So let's list all cells that have letters or ◌.\\ \tt \\ \tt For example, row 1: cells at x=1,2,3,4,5? But x=4 and x=5 are separated by lines, but they are still cells.\\ \tt \\ \tt Actually, the grid has 7 columns and 7 rows of cells, with some walls between.\\ \tt \\ \tt In row 1: the string "C◌◌│◌│◌" means: cell1: C, cell2: ◌, cell3: ◌, then a vertical wall, cell4: ◌, then vertical wall, cell5: ◌? But there are 7 columns, so probably cells 6 and 7 are at the end.\\ \tt \\ \tt Looking at the right edge, it has the border ║ and then the y coordinate.\\ \tt \\ \tt Similarly, row 6: "◌│◌│◌WT" - so cell1: ◌, then vertical wall, cell2: ◌, then vertical wall, cell3: ◌, then cell4: W, cell5: T? But that would be only 5 cells, but we need 7.\\ \tt \\ \tt I think I miscounted.\\ \tt \\ \tt Let's look at the grid carefully:\\ \tt \\ \tt The top border: "╔═══╤═╤═╗" - this indicates that there are vertical lines between columns.\\ \tt \\ \tt The width: from left border to right border, there are 7 characters for the cells and walls.\\ \tt \\ \tt For row 1: "║C◌◌│◌│◌║" - this has 7 characters: position1: C,2:◌,3:◌,4:│,5:◌,6:│,7:◌ but then the right border ║. But the │ are walls, not cells. So the cells are at positions where there is ◌ or a letter.\\ \tt \\ \tt Actually, the cells are the spaces between the walls.\\ \tt \\ \tt The grid has columns 1 to 7, and the walls are between them.\\ \tt \\ \tt In the string, the symbols represent the content of the cell or the wall.\\ \tt \\ \tt For example, in row 1: after ║, the first cell is column 1: 'C', then column 2: '◌', then column 3: '◌', then a wall '│', then column 4: '◌', then wall '│', then column 5: '◌', and then the right border ║. But what about columns 6 and 7? The label says 1 to 7, so there must be 7 columns.\\ \tt \\ \tt Perhaps the "│◌│◌" part covers multiple columns.\\ \tt \\ \tt Another way: the grid is drawn with box-drawing characters, and each "cell" is one character wide.\\ \tt \\ \tt But the cells are not all same width; some have wider spaces.\\ \tt \\ \tt This is confusing.\\ \tt \\ \tt Let's list all cells from the grid:\\ \tt \\ \tt From the left, the first cell is at x=1, and so on.\\ \tt \\ \tt Row 1: after ║, we have "C" which is x=1, then "◌" x=2, then "◌" x=3, then "│" which is a wall between x=3 and x=4, then "◌" which is x=4, then "│" wall between x=4 and x=5, then "◌" which is x=5, and then ║. But x=6 and x=7 are not shown? But the coordinate label says 7 columns.\\ \tt \\ \tt Look at row 4: "║◌◌◌│◌◌◌║" - here we have three ◌, then │, then three ◌, so that's 6 cells? But with the wall, it's 7 characters, but the cells are at positions.\\ \tt \\ \tt Perhaps the walls are drawn in the same space as the cell boundaries.\\ \tt \\ \tt To simplify, I should consider that each cell is represented by one character, and the walls are additional.\\ \tt \\ \tt In row 1: the string has 7 characters between the borders: "C◌◌│◌│◌" which is 6 characters? Let's count: C, ◌, ◌, │, ◌, │, ◌ — that's 7 characters. So:\\ \tt \\ \tt - Char1: cell1 (x=1) : 'C'\\ \tt - Char2: cell2 (x=2) : '◌'\\ \tt - Char3: cell3 (x=3) : '◌'\\ \tt - Char4: wall between x=3 and x=4\\ \tt - Char5: cell4 (x=4) : '◌'\\ \tt - Char6: wall between x=4 and x=5\\ \tt - Char7: cell5 (x=5) : '◌'\\ \tt \\ \tt But the coordinates go up to 7, so where are x=6 and x=7?\\ \tt \\ \tt Now look at row 2: "║◌I◌│◌├─╢" - here, after ║, we have: char1: '◌' (x=1), char2: 'I' (x=2), char3: '◌' (x=3), char4: '│' (wall), char5: '◌' (x=4), char6: '├' (corner), char7: '─' (horizontal wall). So cells only up to x=4? But then the label says y=2, and coordinates to 7.\\ \tt \\ \tt This is not matching.\\ \tt \\ \tt Perhaps the grid has cells only where there is space, and the walls are drawn over the grid.\\ \tt \\ \tt Another idea: the grid has 7 columns, but the drawing uses multiple characters per cell sometimes.\\ \tt \\ \tt For example, in row 6: "║◌│◌│◌WT║" - let's parse: after ║, char1: '◌' (cell x=1), char2: '│' (wall), char3: '◌' (cell x=2), char4: '│' (wall), char5: '◌' (cell x=3), char6: 'W' (cell x=4), char7: 'T' (cell x=5). But then x=6 and x=7 are missing.\\ \tt \\ \tt But the coordinate label on the right says 6, so for row 6, there should be cells x=1 to x=7.\\ \tt \\ \tt I think there is a misunderstanding in the grid representation.\\ \tt \\ \tt Let's look at the full grid:\\ \tt \\ \tt \\ \tt  1234567\\ \tt ╔═══╤═╤═╗\\ \tt ║C◌◌│◌│◌║ 1\\ \tt ║◌I◌│◌├─╢ 2\\ \tt ║◌◌◌├─┘◌║ 3\\ \tt ║◌◌◌│◌◌◌║ 4\\ \tt ╟─┬─┼───╢ 5\\ \tt ║◌│◌│◌WT║ 6\\ \tt ║◌│◌│◌◌H║ 7\\ \tt ╚═╧═╧═══╝\\ \tt \\ \tt \\ \tt The top line " 1234567" indicates the x-coordinates, so there are 7 columns.\\ \tt \\ \tt The second line "╔═══╤═╤═╗" shows the top border. The "═══" might cover multiple columns.\\ \tt \\ \tt In standard box-drawing, "╔" is top-left corner, "═══" is three horizontal lines, "╤" is T-shaped junction, "═" is one horizontal, "╤" another T, "═" one horizontal, and "╗" top-right. So this border has segments of length 3,1,1 respectively, but for 7 columns, it should have 7 segments.\\ \tt \\ \tt Perhaps the grid has cells of varying width, but the coordinates are still 1 to 7.\\ \tt \\ \tt To find the cell positions, we should look at the rows where cells are shown.\\ \tt \\ \tt In row 1: "║C◌◌│◌│◌║ 1" - the "C◌◌" might be three cells in a wide cell, but the coordinate label has separate digits, so likely each digit is one column.\\ \tt \\ \tt I think the best way is to assume that each cell is one unit, and the walls are between them.\\ \tt \\ \tt From row 4: "║◌◌◌│◌◌◌║" - this has three cells, then a wall, then three cells, so that's 6 cells, but we need 7. Notice that the wall "│" is taking the space of one cell, but it is a wall, not a cell.\\ \tt \\ \tt In terms of coordinates, the cells are at positions where there is '◌' or a letter, and the walls are at the boundaries.\\ \tt \\ \tt For example, in row 1, there are cells at x=1,2,3,4,5 because there are five spots with 'C' or '◌', and the walls are at x=3.5 and x=4.5 approximately, but since coordinates are integers, we must have cells only at integer coordinates.\\ \tt \\ \tt This is not working.\\ \tt \\ \tt Let's list the cells from the grid by their content:\\ \tt \\ \tt - At (1,1): 'C'  (from row1, first character after ║ is 'C')\\ \tt - (2,1): '◌' (second character)\\ \tt - (3,1): '◌' (third character)\\ \tt - then '│' is wall, so not a cell\\ \tt - (4,1): '◌' (fifth character? let's index the string.\\ \tt \\ \tt Row 1 string: "║C◌◌│◌│◌║"\\ \tt Positions: 1:║, 2:C, 3:◌, 4:◌, 5:│, 6:◌, 7:│, 8:◌, 9:║\\ \tt \\ \tt But the row has label 1, and x-coordinates 1-7, so likely the cells are at positions 2,3,4,6,8 which correspond to x=1 to x=5, but we need x=6 and x=7.\\ \tt \\ \tt This is not adding up.\\ \tt \\ \tt Perhaps the "═" and "─" are used for horizontal walls, and "│" for vertical, and cells are where they intersect.\\ \tt \\ \tt For example, in row 5: "╟─┬─┼───╢" - this is a border row with junctions.\\ \tt \\ \tt To avoid wasting time, I should use the given grid and assume that the objects are placed in the cells that have letters or dots.\\ \tt \\ \tt From the grid:\\ \tt \\ \tt - 'C' is at (1,1) - first row, first cell\\ \tt - 'I' is at (2,2) - second row, second cell (from row2: "◌I◌..." so at position after ║, second character is 'I')\\ \tt - 'W' and 'T' are at row6, near the end: "◌WT" so at x=6 and x=7 for row6? But in the string, after "│◌│◌" which is cells and walls, then "WT" so 'W' at some x, 'T' at next.\\ \tt \\ \tt Let's assume that for row6, the string "║◌│◌│◌WT║" has:\\ \tt - after ║, first character '◌' is cell (1,6)\\ \tt - then '│' is wall\\ \tt - then '◌' is cell (2,6)\\ \tt - then '│' is wall\\ \tt - then '◌' is cell (3,6)\\ \tt - then 'W' is cell (4,6)\\ \tt - then 'T' is cell (5,6)\\ \tt - but the right border, so no more cells? But label says x=7.\\ \tt \\ \tt Then for row7: "║◌│◌│◌◌H║" - similarly, '◌' at (1,7), wall, '◌' at (2,7), wall, '◌' at (3,7), '◌' at (4,7), 'H' at (5,7).\\ \tt \\ \tt But the coordinates go to 7, so where are x=6 and x=7 for rows 6 and 7?\\ \tt \\ \tt Perhaps the last cells are combined.\\ \tt \\ \tt Notice that in row 4: "║◌◌◌│◌◌◌║" - here, there are 6 ◌, so likely 6 cells, with a wall in the middle.\\ \tt \\ \tt Then for row 4, cells at x=1,2,3,4,5,6? But the wall is at x=4, so cells x=1,2,3 on left, and x=5,6,7 on right? But there are only three on right, so x=5,6,7.\\ \tt \\ \tt Similarly, for row 1: "C◌◌│◌│◌" - here, "C◌◌" for x=1,2,3, then wall, then "◌" for x=4, then wall, then "◌" for x=5, and x=6 and x=7 are not shown, but from the border, they might be empty or something.\\ \tt \\ \tt This is very confusing.\\ \tt \\ \tt Perhaps the grid is not full of cells; some cells are merged.\\ \tt \\ \tt To move objects, we need to know the empty cells.\\ \tt \\ \tt But for the sake of the game, I should use the positions as per the visible letters.\\ \tt \\ \tt From the grid:\\ \tt \\ \tt - C: at (1,1)\\ \tt - I: at (2,2)  (from row2, second character)\\ \tt - W: at (6,6) from row6, but let's see row6: the last part "WT" so W at (6,6) and T at (7,6) if we assume that after the walls, there are cells at x=6 and x=7.\\ \tt \\ \tt Similarly, H at (7,7) from row7: "◌◌H" at the end, so H at (7,7).\\ \tt \\ \tt And T is at (7,6) or (6,6)? In row6, "WT" so likely W at (6,6), T at (7,6).\\ \tt \\ \tt But in row6, there is also "│◌│◌" which are cells at x=1,2,3, so for row6, cells at x=1,2,3,6,7.\\ \tt \\ \tt Similarly for row7, cells at x=1,2,3,6,7? But "◌◌H" might be x=5,6,7? But from the drawing, it's not clear.\\ \tt \\ \tt To resolve this, I should look for the '◌' which indicates empty cell.\\ \tt \\ \tt In row 3: "║◌◌◌├─┘◌║" - here, "◌◌◌" for x=1,2,3, then "├" which is a corner, then "─" horizontal, then "┘" corner, then "◌" for x=4? But then the label.\\ \tt \\ \tt I think I need to accept the grid as is and proceed.\\ \tt \\ \tt For my grid, let's list the object positions based on the letters:\\ \tt \\ \tt - C: at (1,1)\\ \tt - I: at (2,2)  (from row2: "◌I◌..." so at y=2, x=2)\\ \tt - W: at (6,6) from row6: the "W" in "WT" at the end, so assuming x=6,y=6\\ \tt - T: at (7,6) from row6: the "T" in "WT", so x=7,y=6\\ \tt - H: at (7,7) from row7: the "H" in "◌◌H" at the end, so x=7,y=7\\ \tt \\ \tt But are there cells at (4,6) or (5,6)? In row6, after "│◌│◌" which is cells at x=1,2,3, then "W" at x=4? But the string has "◌│◌│◌WT" so after the third "◌", it's "W", so likely x=4 for W, but then "T" at x=5, and no x=6 and x=7.\\ \tt \\ \tt This is not consistent with the coordinate label.\\ \tt \\ \tt Perhaps the coordinate label indicates the column number, and the cells are at those columns, but the drawing might have cells skipped.\\ \tt \\ \tt Another idea: the "═" and "─" represent horizontal spans, but for simplicity, in this game, we should consider only the cells that are explicitly shown with '◌' or letters, and assume that the coordinates are as per the label.\\ \tt \\ \tt So for my grid, let's assume:\\ \tt \\ \tt - C at (1,1)\\ \tt - I at (2,2)\\ \tt - W at (6,6)  // from the "W" in row6\\ \tt - T at (7,6)  // from the "T" in row6\\ \tt - H at (7,7)  // from the "H" in row7\\ \tt \\ \tt But from row6, the "WT" is at the end, so for y=6, x=6 and x=7 have W and T.\\ \tt \\ \tt Similarly, for y=7, x=7 has H.\\ \tt \\ \tt Now for empty cells, we can move objects to any cell that has '◌' in the grid.\\ \tt \\ \tt For example, in row1, there are '◌' at x=2,3,4,5 (assuming from the string).\\ \tt \\ \tt But to avoid errors, I should communicate with the other player based on what I see.\\ \tt \\ \tt The other player has W at (2,2), I at (5,6), T at (1,6), C at (2,1), H at (6,4).\\ \tt \\ \tt So for me, I have C at (1,1), I at (2,2), W at (6,6), T at (7,6), H at (7,7).\\ \tt \\ \tt So let's reply with that.\\ \tt \\ \tt So my positions: C (1,1), I (2,2), W (6,6), T (7,6), H (7,7)\\ \tt \\ \tt Now, we need to agree on common coordinates.\\ \tt \\ \tt Let's choose target positions that are empty on both grids.\\ \tt \\ \tt First, list empty cells on my grid from the drawing:\\ \tt \\ \tt From row1: cells at x=2,3,4,5 are empty? But from string "C◌◌│◌│◌" so after C, ◌ at x=2, ◌ at x=3, then wall, then ◌ at x=4, then wall, then ◌ at x=5. So empty at (2,1), (3,1), (4,1), (5,1)\\ \tt \\ \tt But y=1 has only up to x=5? What about x=6 and x=7? From the border, they might not be cells, or are empty.\\ \tt \\ \tt In row1, after the last ◌, there is ║, so no cells for x=6 and x=7 in row1.\\ \tt \\ \tt Similarly, row2: "◌I◌│◌├─╢" so empty at (1,2), (3,2), (4,2) and then corner, so no more cells.\\ \tt \\ \tt Row3: "◌◌◌├─┘◌║" so empty at (1,3), (2,3), (3,3), and then at the end "◌" which might be (4,3) or something.\\ \tt \\ \tt This is messy.\\ \tt \\ \tt Perhaps for this game, we can use the coordinates from 1 to 7 for x and y, and assume that all cells are available unless occupied by an object or wall, but from the grid, there are walls that block cells.\\ \tt \\ \tt To move, we need to find empty cells that are not walls.\\ \tt \\ \tt From the grid, the cells that have '◌' are empty and can be moved to.\\ \tt \\ \tt So let's list all empty cells from the grid:\\ \tt \\ \tt Row1: at (2,1), (3,1), (4,1), (5,1)  [ since (1,1) has C ]\\ \tt \\ \tt Row2: at (1,2), (3,2), (4,2)  [ (2,2) has I ]\\ \tt \\ \tt Row3: from "◌◌◌" so (1,3), (2,3), (3,3), and then "◌" at the end, so (4,3) or (5,3)? The string has "├─┘◌" so after "├", "─", "┘", then "◌", so likely (4,3) is empty.\\ \tt \\ \tt Row4: "◌◌◌│◌◌◌" so (1,4), (2,4), (3,4), and then (5,4), (6,4), (7,4)?\\ \tt  But the wall is between, so cells at x=1,2,3,5,6,7 for y=4.\\ \tt \\ \tt Row5: "╟─┬─┼───╢" this is a border row, so no cells? Or has cells? The "┬" and "┼" are junctions, so likely no empty cells here, or perhaps cells at the points.\\ \tt \\ \tt But typically, in such grids, row5 might have cells at the intersections.\\ \tt \\ \tt Row6: "◌│◌│◌WT" so empty at (1,6), (2,6), (3,6), and then W and T at (6,6) and (7,6), so no empty at x=4,5 for y=6.\\ \tt \\ \tt Row7: "◌│◌│◌◌H" so empty at (1,7), (2,7), (3,7), (4,7), and H at (7,7), so empty at (5,7)? "◌◌H" so two ◌ and H, so empty at (5,7), (6,7), and H at (7,7).\\ \tt \\ \tt So summarizing empty cells on my grid:\\ \tt y=1: x=2,3,4,5\\ \tt y=2: x=1,3,4\\ \tt y=3: x=1,2,3,4  (assuming (4,3) is empty)\\ \tt y=4: x=1,2,3,5,6,7\\ \tt y=5: no cells or not sure? From row5, it's a border, so probably no cells for y=5.\\ \tt y=6: x=1,2,3\\ \tt y=7: x=1,2,3,5,6\\ \tt \\ \tt And objects at:\\ \tt C (1,1)\\ \tt I (2,2)\\ \tt W (6,6)\\ \tt T (7,6)\\ \tt H (7,7)\\ \tt \\ \tt For the other player, they have objects at:\\ \tt C (2,1)\\ \tt I (5,6)\\ \tt T (1,6)\\ \tt W (2,2)\\ \tt H (6,4)\\ \tt \\ \tt So to agree on common targets, we should choose positions that are empty on both grids.\\ \tt \\ \tt But since we don't know each other's empty cells, we need to choose positions that are likely empty.\\ \tt \\ \tt Common empty cells might be around the center.\\ \tt \\ \tt Let's choose for example (4,4) for all, but (4,4) on my grid is empty? From above, y=4, x=4 is not listed; I have for y=4: x=1,2,3,5,6,7, so x=4 is wall? From row4: "│" is at x=4, so it is a wall, not empty.\\ \tt \\ \tt So (4,4) is not available.\\ \tt \\ \tt Similarly, (3,3) is empty on my grid, and might be on other.\\ \tt \\ \tt But let's first reply with my positions.\\ \tt \\ \tt So I'll say my positions.         ' 
	  } 
	   } 
	   } 
	 & & \\ 
 

    \theutterance \stepcounter{utterance}  

    & & \multicolumn{4}{p{0.6\linewidth}}{\cellcolor[rgb]{0.9,0.9,0.9}{%
	\makecell[{{p{\linewidth}}}]{% 
	  \tt {\tiny [A$\langle$GM]}  
	 Penalty: Your message must not contain anything before the command!\\ \tt \\ \tt Make sure that your response only contains either `SAY: <MESSAGE>` or `MOVE: <OBJECT>, (<X>, <Y>)`, and nothing else!\\ \tt \\ \tt You have collectively accumulated 2 of 12 penalties.\\ \tt Please try again! 
	  } 
	   } 
	   } 
	 & \\ 
 

    \theutterance \stepcounter{utterance}  

    & \multicolumn{4}{p{0.6\linewidth}}{\cellcolor[rgb]{0.8,1,0.9}{%
	\makecell[{{p{\linewidth}}}]{% 
	  \tt {\tiny [A$\rangle$GM]}  
	 SAY: My object positions: C at (1,1), I at (2,2), W at (6,6), T at (7,6), H at (7,7). What are your empty cells? We can choose target positions that are empty for both. For example, (3,3) might be good for some. 
	  } 
	   } 
	   } 
	 & & \\ 
 

    \theutterance \stepcounter{utterance}  

    & & & \multicolumn{2}{p{0.3\linewidth}}{\cellcolor[rgb]{0.95,0.95,0.95}{%
	\makecell[{{p{\linewidth}}}]{% 
	  \tt {\tiny [GM$|$GM]}  
	 SAY: My object positions: C at (1,1), I at (2,2), W at (6,6), T at (7,6), H at (7,7). What are your empty cells? We can choose target positions that are empty for both. For example, (3,3) might be good for some. 
	  } 
	   } 
	   } 
	 & & \\ 
 

    \theutterance \stepcounter{utterance}  

    & & \multicolumn{4}{p{0.6\linewidth}}{\cellcolor[rgb]{0.9,0.9,0.9}{%
	\makecell[{{p{\linewidth}}}]{% 
	  \tt {\tiny [GM$\rangle$B]}  
	 Your message has been relayed to the other player.\\ \tt \\ \tt \\ \tt You are currently playing round 2 of maximum 20.\\ \tt You have collectively accumulated 2 of 12 penalties.\\ \tt The other player sent this message:\\ \tt \\ \tt \\ \tt "My object positions: C at (1,1), I at (2,2), W at (6,6), T at (7,6), H at (7,7). What are your empty cells? We can choose target positions that are empty for both. For example, (3,3) might be good for some."\\ \tt \\ \tt \\ \tt What is your next command? 
	  } 
	   } 
	   } 
	 & \\ 
 

    \theutterance \stepcounter{utterance}  

    & & & \multicolumn{4}{p{0.6\linewidth}}{\cellcolor[rgb]{1,0.85,0.72}{%
	\makecell[{{p{\linewidth}}}]{% 
	  \tt {\tiny [GM$\langle$B]}  
	 SAY: My empty cells are all positions without objects. For targets, let's use a grid: (1,1) to (7,7). I propose: C to (1,1), W to (2,2), I to (3,3), T to (4,4), H to (5,5). Confirm or suggest alternatives? 
	  } 
	   } 
	   } 
	  \\ 
 

    \theutterance \stepcounter{utterance}  

    & & & \multicolumn{2}{p{0.3\linewidth}}{\cellcolor[rgb]{0.95,0.95,0.95}{%
	\makecell[{{p{\linewidth}}}]{% 
	  \tt {\tiny [GM$|$GM]}  
	 SAY: My empty cells are all positions without objects. For targets, let's use a grid: (1,1) to (7,7). I propose: C to (1,1), W to (2,2), I to (3,3), T to (4,4), H to (5,5). Confirm or suggest alternatives? 
	  } 
	   } 
	   } 
	 & & \\ 
 

    \theutterance \stepcounter{utterance}  

    & & \multicolumn{4}{p{0.6\linewidth}}{\cellcolor[rgb]{0.9,0.9,0.9}{%
	\makecell[{{p{\linewidth}}}]{% 
	  \tt {\tiny [A$\langle$GM]}  
	 Your message has been relayed to the other player.\\ \tt \\ \tt \\ \tt You are currently playing round 2 of maximum 20.\\ \tt You have collectively accumulated 2 of 12 penalties.\\ \tt The other player sent this message:\\ \tt \\ \tt \\ \tt "My empty cells are all positions without objects. For targets, let's use a grid: (1,1) to (7,7). I propose: C to (1,1), W to (2,2), I to (3,3), T to (4,4), H to (5,5). Confirm or suggest alternatives?"\\ \tt \\ \tt \\ \tt What is your next command? 
	  } 
	   } 
	   } 
	 & \\ 
 

    \theutterance \stepcounter{utterance}  

    & \multicolumn{4}{p{0.6\linewidth}}{\cellcolor[rgb]{0.8,1,0.9}{%
	\makecell[{{p{\linewidth}}}]{% 
	  \tt {\tiny [A$\rangle$GM]}  
	 SAY: Confirmed. I will move objects to those targets. I start with moving C to (1,1) but it's already there. Then I move W to (2,2) but it's occupied by I. I need to move I first. Let's sequence moves to avoid conflicts. I'll move I to (3,3) first. 
	  } 
	   } 
	   } 
	 & & \\ 
 

    \theutterance \stepcounter{utterance}  

    & & & \multicolumn{2}{p{0.3\linewidth}}{\cellcolor[rgb]{0.95,0.95,0.95}{%
	\makecell[{{p{\linewidth}}}]{% 
	  \tt {\tiny [GM$|$GM]}  
	 SAY: Confirmed. I will move objects to those targets. I start with moving C to (1,1) but it's already there. Then I move W to (2,2) but it's occupied by I. I need to move I first. Let's sequence moves to avoid conflicts. I'll move I to (3,3) first. 
	  } 
	   } 
	   } 
	 & & \\ 
 

    \theutterance \stepcounter{utterance}  

    & & \multicolumn{4}{p{0.6\linewidth}}{\cellcolor[rgb]{0.9,0.9,0.9}{%
	\makecell[{{p{\linewidth}}}]{% 
	  \tt {\tiny [GM$\rangle$B]}  
	 Your message has been relayed to the other player.\\ \tt \\ \tt \\ \tt You are currently playing round 3 of maximum 20.\\ \tt You have collectively accumulated 2 of 12 penalties.\\ \tt The other player sent this message:\\ \tt \\ \tt \\ \tt "Confirmed. I will move objects to those targets. I start with moving C to (1,1) but it's already there. Then I move W to (2,2) but it's occupied by I. I need to move I first. Let's sequence moves to avoid conflicts. I'll move I to (3,3) first."\\ \tt \\ \tt \\ \tt What is your next command? 
	  } 
	   } 
	   } 
	 & \\ 
 

    \theutterance \stepcounter{utterance}  

    & & & \multicolumn{4}{p{0.6\linewidth}}{\cellcolor[rgb]{1,0.85,0.72}{%
	\makecell[{{p{\linewidth}}}]{% 
	  \tt {\tiny [GM$\langle$B]}  
	 SAY: Agreed. I will also sequence moves to avoid conflicts. I'll start by moving my I from (5,6) to (3,3). After that, I'll move my W from (2,2) to (2,2) is already correct. Then others. 
	  } 
	   } 
	   } 
	  \\ 
 

    \theutterance \stepcounter{utterance}  

    & & & \multicolumn{2}{p{0.3\linewidth}}{\cellcolor[rgb]{0.95,0.95,0.95}{%
	\makecell[{{p{\linewidth}}}]{% 
	  \tt {\tiny [GM$|$GM]}  
	 SAY: Agreed. I will also sequence moves to avoid conflicts. I'll start by moving my I from (5,6) to (3,3). After that, I'll move my W from (2,2) to (2,2) is already correct. Then others. 
	  } 
	   } 
	   } 
	 & & \\ 
 

    \theutterance \stepcounter{utterance}  

    & & \multicolumn{4}{p{0.6\linewidth}}{\cellcolor[rgb]{0.9,0.9,0.9}{%
	\makecell[{{p{\linewidth}}}]{% 
	  \tt {\tiny [A$\langle$GM]}  
	 Your message has been relayed to the other player.\\ \tt \\ \tt \\ \tt You are currently playing round 3 of maximum 20.\\ \tt You have collectively accumulated 2 of 12 penalties.\\ \tt The other player sent this message:\\ \tt \\ \tt \\ \tt "Agreed. I will also sequence moves to avoid conflicts. I'll start by moving my I from (5,6) to (3,3). After that, I'll move my W from (2,2) to (2,2) is already correct. Then others."\\ \tt \\ \tt \\ \tt What is your next command? 
	  } 
	   } 
	   } 
	 & \\ 
 

    \theutterance \stepcounter{utterance}  

    & \multicolumn{4}{p{0.6\linewidth}}{\cellcolor[rgb]{0.8,1,0.9}{%
	\makecell[{{p{\linewidth}}}]{% 
	  \tt {\tiny [A$\rangle$GM]}  
	 MOVE: I, (3, 3) 
	  } 
	   } 
	   } 
	 & & \\ 
 

    \theutterance \stepcounter{utterance}  

    & & & \multicolumn{2}{p{0.3\linewidth}}{\cellcolor[rgb]{0.95,0.95,0.95}{%
	\makecell[{{p{\linewidth}}}]{% 
	  \tt {\tiny [GM$|$GM]}  
	 MOVE: I, (3, 3) 
	  } 
	   } 
	   } 
	 & & \\ 
 

    \theutterance \stepcounter{utterance}  

    & & & \multicolumn{2}{p{0.3\linewidth}}{\cellcolor[rgb]{0.95,0.95,0.95}{%
	\makecell[{{p{\linewidth}}}]{% 
	  \tt {\tiny [GM$|$GM]}  
	 Moved 'I' to (3,3) successfully. Your updated grid looks like this:\\ \tt ```\\ \tt  1234567\\ \tt ╔═══╤═╤═╗\\ \tt ║C◌◌│◌│◌║ 1\\ \tt ║◌◌◌│◌├─╢ 2\\ \tt ║◌◌I├─┘◌║ 3\\ \tt ║◌◌◌│◌◌◌║ 4\\ \tt ╟─┬─┼───╢ 5\\ \tt ║◌│◌│◌WT║ 6\\ \tt ║◌│◌│◌◌H║ 7\\ \tt ╚═╧═╧═══╝\\ \tt ```\\ \tt  
	  } 
	   } 
	   } 
	 & & \\ 
 

    \theutterance \stepcounter{utterance}  

    & & \multicolumn{4}{p{0.6\linewidth}}{\cellcolor[rgb]{0.9,0.9,0.9}{%
	\makecell[{{p{\linewidth}}}]{% 
	  \tt {\tiny [GM$\rangle$B]}  
	 Your message has been relayed to the other player.\\ \tt \\ \tt \\ \tt You are currently playing round 4 of maximum 20.\\ \tt You have collectively accumulated 2 of 12 penalties.\\ \tt The other player moved an object on their grid.\\ \tt \\ \tt \\ \tt What is your next command? 
	  } 
	   } 
	   } 
	 & \\ 
 

    \theutterance \stepcounter{utterance}  

    & & & \multicolumn{4}{p{0.6\linewidth}}{\cellcolor[rgb]{1,0.85,0.72}{%
	\makecell[{{p{\linewidth}}}]{% 
	  \tt {\tiny [GM$\langle$B]}  
	 MOVE: I, (3,3) 
	  } 
	   } 
	   } 
	  \\ 
 

    \theutterance \stepcounter{utterance}  

    & & & \multicolumn{2}{p{0.3\linewidth}}{\cellcolor[rgb]{0.95,0.95,0.95}{%
	\makecell[{{p{\linewidth}}}]{% 
	  \tt {\tiny [GM$|$GM]}  
	 MOVE: I, (3,3) 
	  } 
	   } 
	   } 
	 & & \\ 
 

    \theutterance \stepcounter{utterance}  

    & & & \multicolumn{2}{p{0.3\linewidth}}{\cellcolor[rgb]{0.95,0.95,0.95}{%
	\makecell[{{p{\linewidth}}}]{% 
	  \tt {\tiny [GM$|$GM]}  
	 Moved 'I' to (3,3) successfully. Your updated grid looks like this:\\ \tt ```\\ \tt  1234567\\ \tt ╔═══╤═╤═╗\\ \tt ║◌◌◌│C│◌║ 1\\ \tt ║◌W◌│◌├─╢ 2\\ \tt ║◌◌I├─┘◌║ 3\\ \tt ║◌◌◌│◌◌H║ 4\\ \tt ╟─┬─┼───╢ 5\\ \tt ║T│◌│◌◌◌║ 6\\ \tt ║◌│◌│◌◌◌║ 7\\ \tt ╚═╧═╧═══╝\\ \tt ```\\ \tt  
	  } 
	   } 
	   } 
	 & & \\ 
 

    \theutterance \stepcounter{utterance}  

    & & \multicolumn{4}{p{0.6\linewidth}}{\cellcolor[rgb]{0.9,0.9,0.9}{%
	\makecell[{{p{\linewidth}}}]{% 
	  \tt {\tiny [A$\langle$GM]}  
	 Moved 'I' to (3,3) successfully. Your updated grid looks like this:\\ \tt ```\\ \tt  1234567\\ \tt ╔═══╤═╤═╗\\ \tt ║C◌◌│◌│◌║ 1\\ \tt ║◌◌◌│◌├─╢ 2\\ \tt ║◌◌I├─┘◌║ 3\\ \tt ║◌◌◌│◌◌◌║ 4\\ \tt ╟─┬─┼───╢ 5\\ \tt ║◌│◌│◌WT║ 6\\ \tt ║◌│◌│◌◌H║ 7\\ \tt ╚═╧═╧═══╝\\ \tt ```\\ \tt You are currently playing round 4 of maximum 20.\\ \tt You have collectively accumulated 2 of 12 penalties.\\ \tt The other player moved an object on their grid.\\ \tt \\ \tt \\ \tt What is your next command? 
	  } 
	   } 
	   } 
	 & \\ 
 

    \theutterance \stepcounter{utterance}  

    & \multicolumn{4}{p{0.6\linewidth}}{\cellcolor[rgb]{0.8,1,0.9}{%
	\makecell[{{p{\linewidth}}}]{% 
	  \tt {\tiny [A$\rangle$GM]}  
	 MOVE: W, (2, 2) 
	  } 
	   } 
	   } 
	 & & \\ 
 

    \theutterance \stepcounter{utterance}  

    & & & \multicolumn{2}{p{0.3\linewidth}}{\cellcolor[rgb]{0.95,0.95,0.95}{%
	\makecell[{{p{\linewidth}}}]{% 
	  \tt {\tiny [GM$|$GM]}  
	 MOVE: W, (2, 2) 
	  } 
	   } 
	   } 
	 & & \\ 
 

    \theutterance \stepcounter{utterance}  

    & & & \multicolumn{2}{p{0.3\linewidth}}{\cellcolor[rgb]{0.95,0.95,0.95}{%
	\makecell[{{p{\linewidth}}}]{% 
	  \tt {\tiny [GM$|$GM]}  
	 Moved 'W' to (2,2) successfully. Your updated grid looks like this:\\ \tt ```\\ \tt  1234567\\ \tt ╔═══╤═╤═╗\\ \tt ║C◌◌│◌│◌║ 1\\ \tt ║◌W◌│◌├─╢ 2\\ \tt ║◌◌I├─┘◌║ 3\\ \tt ║◌◌◌│◌◌◌║ 4\\ \tt ╟─┬─┼───╢ 5\\ \tt ║◌│◌│◌◌T║ 6\\ \tt ║◌│◌│◌◌H║ 7\\ \tt ╚═╧═╧═══╝\\ \tt ```\\ \tt  
	  } 
	   } 
	   } 
	 & & \\ 
 

    \theutterance \stepcounter{utterance}  

    & & \multicolumn{4}{p{0.6\linewidth}}{\cellcolor[rgb]{0.9,0.9,0.9}{%
	\makecell[{{p{\linewidth}}}]{% 
	  \tt {\tiny [GM$\rangle$B]}  
	 Moved 'I' to (3,3) successfully. Your updated grid looks like this:\\ \tt ```\\ \tt  1234567\\ \tt ╔═══╤═╤═╗\\ \tt ║◌◌◌│C│◌║ 1\\ \tt ║◌W◌│◌├─╢ 2\\ \tt ║◌◌I├─┘◌║ 3\\ \tt ║◌◌◌│◌◌H║ 4\\ \tt ╟─┬─┼───╢ 5\\ \tt ║T│◌│◌◌◌║ 6\\ \tt ║◌│◌│◌◌◌║ 7\\ \tt ╚═╧═╧═══╝\\ \tt ```\\ \tt You are currently playing round 5 of maximum 20.\\ \tt You have collectively accumulated 2 of 12 penalties.\\ \tt The other player moved an object on their grid.\\ \tt \\ \tt \\ \tt What is your next command? 
	  } 
	   } 
	   } 
	 & \\ 
 

    \theutterance \stepcounter{utterance}  

    & & & \multicolumn{4}{p{0.6\linewidth}}{\cellcolor[rgb]{1,0.85,0.72}{%
	\makecell[{{p{\linewidth}}}]{% 
	  \tt {\tiny [GM$\langle$B]}  
	 SAY: I moved I to (3,3). My W is already at (2,2). Next, I'll move T to (4,4). What is your status? 
	  } 
	   } 
	   } 
	  \\ 
 

    \theutterance \stepcounter{utterance}  

    & & & \multicolumn{2}{p{0.3\linewidth}}{\cellcolor[rgb]{0.95,0.95,0.95}{%
	\makecell[{{p{\linewidth}}}]{% 
	  \tt {\tiny [GM$|$GM]}  
	 SAY: I moved I to (3,3). My W is already at (2,2). Next, I'll move T to (4,4). What is your status? 
	  } 
	   } 
	   } 
	 & & \\ 
 

    \theutterance \stepcounter{utterance}  

    & & \multicolumn{4}{p{0.6\linewidth}}{\cellcolor[rgb]{0.9,0.9,0.9}{%
	\makecell[{{p{\linewidth}}}]{% 
	  \tt {\tiny [A$\langle$GM]}  
	 Moved 'W' to (2,2) successfully. Your updated grid looks like this:\\ \tt ```\\ \tt  1234567\\ \tt ╔═══╤═╤═╗\\ \tt ║C◌◌│◌│◌║ 1\\ \tt ║◌W◌│◌├─╢ 2\\ \tt ║◌◌I├─┘◌║ 3\\ \tt ║◌◌◌│◌◌◌║ 4\\ \tt ╟─┬─┼───╢ 5\\ \tt ║◌│◌│◌◌T║ 6\\ \tt ║◌│◌│◌◌H║ 7\\ \tt ╚═╧═╧═══╝\\ \tt ```\\ \tt You are currently playing round 5 of maximum 20.\\ \tt You have collectively accumulated 2 of 12 penalties.\\ \tt The other player sent this message:\\ \tt \\ \tt \\ \tt "I moved I to (3,3). My W is already at (2,2). Next, I'll move T to (4,4). What is your status?"\\ \tt \\ \tt \\ \tt What is your next command? 
	  } 
	   } 
	   } 
	 & \\ 
 

    \theutterance \stepcounter{utterance}  

    & \multicolumn{4}{p{0.6\linewidth}}{\cellcolor[rgb]{0.8,1,0.9}{%
	\makecell[{{p{\linewidth}}}]{% 
	  \tt {\tiny [A$\rangle$GM]}  
	 MOVE: T, (4, 4) 
	  } 
	   } 
	   } 
	 & & \\ 
 

    \theutterance \stepcounter{utterance}  

    & & & \multicolumn{2}{p{0.3\linewidth}}{\cellcolor[rgb]{0.95,0.95,0.95}{%
	\makecell[{{p{\linewidth}}}]{% 
	  \tt {\tiny [GM$|$GM]}  
	 MOVE: T, (4, 4) 
	  } 
	   } 
	   } 
	 & & \\ 
 

    \theutterance \stepcounter{utterance}  

    & & & \multicolumn{2}{p{0.3\linewidth}}{\cellcolor[rgb]{0.95,0.95,0.95}{%
	\makecell[{{p{\linewidth}}}]{% 
	  \tt {\tiny [GM$|$GM]}  
	 Penalty: (4,4) is not empty, but contains '│'.\\ \tt \\ \tt You have collectively accumulated 3 of 12 penalties.\\ \tt Please try again! 
	  } 
	   } 
	   } 
	 & & \\ 
 

    \theutterance \stepcounter{utterance}  

    & & \multicolumn{4}{p{0.6\linewidth}}{\cellcolor[rgb]{0.9,0.9,0.9}{%
	\makecell[{{p{\linewidth}}}]{% 
	  \tt {\tiny [A$\langle$GM]}  
	 Penalty: (4,4) is not empty, but contains '│'.\\ \tt \\ \tt You have collectively accumulated 3 of 12 penalties.\\ \tt Please try again! 
	  } 
	   } 
	   } 
	 & \\ 
 

    \theutterance \stepcounter{utterance}  

    & \multicolumn{4}{p{0.6\linewidth}}{\cellcolor[rgb]{0.8,1,0.9}{%
	\makecell[{{p{\linewidth}}}]{% 
	  \tt {\tiny [A$\rangle$GM]}  
	 SAY: (4,4) is a wall on my grid. Let's choose a different target for T. How about (5,5)? 
	  } 
	   } 
	   } 
	 & & \\ 
 

    \theutterance \stepcounter{utterance}  

    & & & \multicolumn{2}{p{0.3\linewidth}}{\cellcolor[rgb]{0.95,0.95,0.95}{%
	\makecell[{{p{\linewidth}}}]{% 
	  \tt {\tiny [GM$|$GM]}  
	 SAY: (4,4) is a wall on my grid. Let's choose a different target for T. How about (5,5)? 
	  } 
	   } 
	   } 
	 & & \\ 
 

    \theutterance \stepcounter{utterance}  

    & & \multicolumn{4}{p{0.6\linewidth}}{\cellcolor[rgb]{0.9,0.9,0.9}{%
	\makecell[{{p{\linewidth}}}]{% 
	  \tt {\tiny [GM$\rangle$B]}  
	 Your message has been relayed to the other player.\\ \tt \\ \tt \\ \tt You are currently playing round 6 of maximum 20.\\ \tt You have collectively accumulated 3 of 12 penalties.\\ \tt The other player sent this message:\\ \tt \\ \tt \\ \tt "(4,4) is a wall on my grid. Let's choose a different target for T. How about (5,5)?"\\ \tt \\ \tt \\ \tt What is your next command? 
	  } 
	   } 
	   } 
	 & \\ 
 

    \theutterance \stepcounter{utterance}  

    & & & \multicolumn{4}{p{0.6\linewidth}}{\cellcolor[rgb]{1,0.85,0.72}{%
	\makecell[{{p{\linewidth}}}]{% 
	  \tt {\tiny [GM$\langle$B]}  
	 SAY: (5,5) is a wall on my grid (it's a horizontal line). Let's use (6,6) for T instead. I'll move T to (6,6). Confirm? 
	  } 
	   } 
	   } 
	  \\ 
 

    \theutterance \stepcounter{utterance}  

    & & & \multicolumn{2}{p{0.3\linewidth}}{\cellcolor[rgb]{0.95,0.95,0.95}{%
	\makecell[{{p{\linewidth}}}]{% 
	  \tt {\tiny [GM$|$GM]}  
	 SAY: (5,5) is a wall on my grid (it's a horizontal line). Let's use (6,6) for T instead. I'll move T to (6,6). Confirm? 
	  } 
	   } 
	   } 
	 & & \\ 
 

    \theutterance \stepcounter{utterance}  

    & & \multicolumn{4}{p{0.6\linewidth}}{\cellcolor[rgb]{0.9,0.9,0.9}{%
	\makecell[{{p{\linewidth}}}]{% 
	  \tt {\tiny [A$\langle$GM]}  
	 Your message has been relayed to the other player.\\ \tt \\ \tt \\ \tt You are currently playing round 6 of maximum 20.\\ \tt You have collectively accumulated 3 of 12 penalties.\\ \tt The other player sent this message:\\ \tt \\ \tt \\ \tt "(5,5) is a wall on my grid (it's a horizontal line). Let's use (6,6) for T instead. I'll move T to (6,6). Confirm?"\\ \tt \\ \tt \\ \tt What is your next command? 
	  } 
	   } 
	   } 
	 & \\ 
 

    \theutterance \stepcounter{utterance}  

    & \multicolumn{4}{p{0.6\linewidth}}{\cellcolor[rgb]{0.8,1,0.9}{%
	\makecell[{{p{\linewidth}}}]{% 
	  \tt {\tiny [A$\rangle$GM]}  
	 SAY: Confirmed. T to (6,6). I will move T to (6,6) after you confirm it's empty for you. 
	  } 
	   } 
	   } 
	 & & \\ 
 

    \theutterance \stepcounter{utterance}  

    & & & \multicolumn{2}{p{0.3\linewidth}}{\cellcolor[rgb]{0.95,0.95,0.95}{%
	\makecell[{{p{\linewidth}}}]{% 
	  \tt {\tiny [GM$|$GM]}  
	 SAY: Confirmed. T to (6,6). I will move T to (6,6) after you confirm it's empty for you. 
	  } 
	   } 
	   } 
	 & & \\ 
 

    \theutterance \stepcounter{utterance}  

    & & \multicolumn{4}{p{0.6\linewidth}}{\cellcolor[rgb]{0.9,0.9,0.9}{%
	\makecell[{{p{\linewidth}}}]{% 
	  \tt {\tiny [GM$\rangle$B]}  
	 Your message has been relayed to the other player.\\ \tt \\ \tt \\ \tt You are currently playing round 7 of maximum 20.\\ \tt You have collectively accumulated 3 of 12 penalties.\\ \tt The other player sent this message:\\ \tt \\ \tt \\ \tt "Confirmed. T to (6,6). I will move T to (6,6) after you confirm it's empty for you."\\ \tt \\ \tt \\ \tt What is your next command? 
	  } 
	   } 
	   } 
	 & \\ 
 

    \theutterance \stepcounter{utterance}  

    & & & \multicolumn{4}{p{0.6\linewidth}}{\cellcolor[rgb]{1,0.85,0.72}{%
	\makecell[{{p{\linewidth}}}]{% 
	  \tt {\tiny [GM$\langle$B]}  
	 SAY: (6,6) is empty on my grid. I will move T to (6,6) now. 
	  } 
	   } 
	   } 
	  \\ 
 

    \theutterance \stepcounter{utterance}  

    & & & \multicolumn{2}{p{0.3\linewidth}}{\cellcolor[rgb]{0.95,0.95,0.95}{%
	\makecell[{{p{\linewidth}}}]{% 
	  \tt {\tiny [GM$|$GM]}  
	 SAY: (6,6) is empty on my grid. I will move T to (6,6) now. 
	  } 
	   } 
	   } 
	 & & \\ 
 

    \theutterance \stepcounter{utterance}  

    & & \multicolumn{4}{p{0.6\linewidth}}{\cellcolor[rgb]{0.9,0.9,0.9}{%
	\makecell[{{p{\linewidth}}}]{% 
	  \tt {\tiny [A$\langle$GM]}  
	 Your message has been relayed to the other player.\\ \tt \\ \tt \\ \tt You are currently playing round 7 of maximum 20.\\ \tt You have collectively accumulated 3 of 12 penalties.\\ \tt The other player sent this message:\\ \tt \\ \tt \\ \tt "(6,6) is empty on my grid. I will move T to (6,6) now."\\ \tt \\ \tt \\ \tt What is your next command? 
	  } 
	   } 
	   } 
	 & \\ 
 

    \theutterance \stepcounter{utterance}  

    & \multicolumn{4}{p{0.6\linewidth}}{\cellcolor[rgb]{0.8,1,0.9}{%
	\makecell[{{p{\linewidth}}}]{% 
	  \tt {\tiny [A$\rangle$GM]}  
	 MOVE: T, (6, 6) 
	  } 
	   } 
	   } 
	 & & \\ 
 

    \theutterance \stepcounter{utterance}  

    & & & \multicolumn{2}{p{0.3\linewidth}}{\cellcolor[rgb]{0.95,0.95,0.95}{%
	\makecell[{{p{\linewidth}}}]{% 
	  \tt {\tiny [GM$|$GM]}  
	 MOVE: T, (6, 6) 
	  } 
	   } 
	   } 
	 & & \\ 
 

    \theutterance \stepcounter{utterance}  

    & & & \multicolumn{2}{p{0.3\linewidth}}{\cellcolor[rgb]{0.95,0.95,0.95}{%
	\makecell[{{p{\linewidth}}}]{% 
	  \tt {\tiny [GM$|$GM]}  
	 Moved 'T' to (6,6) successfully. Your updated grid looks like this:\\ \tt ```\\ \tt  1234567\\ \tt ╔═══╤═╤═╗\\ \tt ║C◌◌│◌│◌║ 1\\ \tt ║◌W◌│◌├─╢ 2\\ \tt ║◌◌I├─┘◌║ 3\\ \tt ║◌◌◌│◌◌◌║ 4\\ \tt ╟─┬─┼───╢ 5\\ \tt ║◌│◌│◌T◌║ 6\\ \tt ║◌│◌│◌◌H║ 7\\ \tt ╚═╧═╧═══╝\\ \tt ```\\ \tt  
	  } 
	   } 
	   } 
	 & & \\ 
 

    \theutterance \stepcounter{utterance}  

    & & \multicolumn{4}{p{0.6\linewidth}}{\cellcolor[rgb]{0.9,0.9,0.9}{%
	\makecell[{{p{\linewidth}}}]{% 
	  \tt {\tiny [GM$\rangle$B]}  
	 Your message has been relayed to the other player.\\ \tt \\ \tt \\ \tt You are currently playing round 8 of maximum 20.\\ \tt You have collectively accumulated 3 of 12 penalties.\\ \tt The other player moved an object on their grid.\\ \tt \\ \tt \\ \tt What is your next command? 
	  } 
	   } 
	   } 
	 & \\ 
 

    \theutterance \stepcounter{utterance}  

    & & & \multicolumn{4}{p{0.6\linewidth}}{\cellcolor[rgb]{1,0.85,0.72}{%
	\makecell[{{p{\linewidth}}}]{% 
	  \tt {\tiny [GM$\langle$B]}  
	 MOVE: T, (6,6) 
	  } 
	   } 
	   } 
	  \\ 
 

    \theutterance \stepcounter{utterance}  

    & & & \multicolumn{2}{p{0.3\linewidth}}{\cellcolor[rgb]{0.95,0.95,0.95}{%
	\makecell[{{p{\linewidth}}}]{% 
	  \tt {\tiny [GM$|$GM]}  
	 MOVE: T, (6,6) 
	  } 
	   } 
	   } 
	 & & \\ 
 

    \theutterance \stepcounter{utterance}  

    & & & \multicolumn{2}{p{0.3\linewidth}}{\cellcolor[rgb]{0.95,0.95,0.95}{%
	\makecell[{{p{\linewidth}}}]{% 
	  \tt {\tiny [GM$|$GM]}  
	 Moved 'T' to (6,6) successfully. Your updated grid looks like this:\\ \tt ```\\ \tt  1234567\\ \tt ╔═══╤═╤═╗\\ \tt ║◌◌◌│C│◌║ 1\\ \tt ║◌W◌│◌├─╢ 2\\ \tt ║◌◌I├─┘◌║ 3\\ \tt ║◌◌◌│◌◌H║ 4\\ \tt ╟─┬─┼───╢ 5\\ \tt ║◌│◌│◌T◌║ 6\\ \tt ║◌│◌│◌◌◌║ 7\\ \tt ╚═╧═╧═══╝\\ \tt ```\\ \tt  
	  } 
	   } 
	   } 
	 & & \\ 
 

    \theutterance \stepcounter{utterance}  

    & & \multicolumn{4}{p{0.6\linewidth}}{\cellcolor[rgb]{0.9,0.9,0.9}{%
	\makecell[{{p{\linewidth}}}]{% 
	  \tt {\tiny [A$\langle$GM]}  
	 Moved 'T' to (6,6) successfully. Your updated grid looks like this:\\ \tt ```\\ \tt  1234567\\ \tt ╔═══╤═╤═╗\\ \tt ║C◌◌│◌│◌║ 1\\ \tt ║◌W◌│◌├─╢ 2\\ \tt ║◌◌I├─┘◌║ 3\\ \tt ║◌◌◌│◌◌◌║ 4\\ \tt ╟─┬─┼───╢ 5\\ \tt ║◌│◌│◌T◌║ 6\\ \tt ║◌│◌│◌◌H║ 7\\ \tt ╚═╧═╧═══╝\\ \tt ```\\ \tt You are currently playing round 8 of maximum 20.\\ \tt You have collectively accumulated 3 of 12 penalties.\\ \tt The other player moved an object on their grid.\\ \tt \\ \tt \\ \tt What is your next command? 
	  } 
	   } 
	   } 
	 & \\ 
 

    \theutterance \stepcounter{utterance}  

    & \multicolumn{4}{p{0.6\linewidth}}{\cellcolor[rgb]{0.8,1,0.9}{%
	\makecell[{{p{\linewidth}}}]{% 
	  \tt {\tiny [A$\rangle$GM]}  
	 SAY: T moved to (6,6). Now, for H, let's use (7,7). Is (7,7) empty for you? 
	  } 
	   } 
	   } 
	 & & \\ 
 

    \theutterance \stepcounter{utterance}  

    & & & \multicolumn{2}{p{0.3\linewidth}}{\cellcolor[rgb]{0.95,0.95,0.95}{%
	\makecell[{{p{\linewidth}}}]{% 
	  \tt {\tiny [GM$|$GM]}  
	 SAY: T moved to (6,6). Now, for H, let's use (7,7). Is (7,7) empty for you? 
	  } 
	   } 
	   } 
	 & & \\ 
 

    \theutterance \stepcounter{utterance}  

    & & \multicolumn{4}{p{0.6\linewidth}}{\cellcolor[rgb]{0.9,0.9,0.9}{%
	\makecell[{{p{\linewidth}}}]{% 
	  \tt {\tiny [GM$\rangle$B]}  
	 Moved 'T' to (6,6) successfully. Your updated grid looks like this:\\ \tt ```\\ \tt  1234567\\ \tt ╔═══╤═╤═╗\\ \tt ║◌◌◌│C│◌║ 1\\ \tt ║◌W◌│◌├─╢ 2\\ \tt ║◌◌I├─┘◌║ 3\\ \tt ║◌◌◌│◌◌H║ 4\\ \tt ╟─┬─┼───╢ 5\\ \tt ║◌│◌│◌T◌║ 6\\ \tt ║◌│◌│◌◌◌║ 7\\ \tt ╚═╧═╧═══╝\\ \tt ```\\ \tt You are currently playing round 9 of maximum 20.\\ \tt You have collectively accumulated 3 of 12 penalties.\\ \tt The other player sent this message:\\ \tt \\ \tt \\ \tt "T moved to (6,6). Now, for H, let's use (7,7). Is (7,7) empty for you?"\\ \tt \\ \tt \\ \tt What is your next command? 
	  } 
	   } 
	   } 
	 & \\ 
 

    \theutterance \stepcounter{utterance}  

    & & & \multicolumn{4}{p{0.6\linewidth}}{\cellcolor[rgb]{1,0.85,0.72}{%
	\makecell[{{p{\linewidth}}}]{% 
	  \tt {\tiny [GM$\langle$B]}  
	 SAY: (7,7) is empty for me. I will move H to (7,7). Confirm? 
	  } 
	   } 
	   } 
	  \\ 
 

    \theutterance \stepcounter{utterance}  

    & & & \multicolumn{2}{p{0.3\linewidth}}{\cellcolor[rgb]{0.95,0.95,0.95}{%
	\makecell[{{p{\linewidth}}}]{% 
	  \tt {\tiny [GM$|$GM]}  
	 SAY: (7,7) is empty for me. I will move H to (7,7). Confirm? 
	  } 
	   } 
	   } 
	 & & \\ 
 

    \theutterance \stepcounter{utterance}  

    & & \multicolumn{4}{p{0.6\linewidth}}{\cellcolor[rgb]{0.9,0.9,0.9}{%
	\makecell[{{p{\linewidth}}}]{% 
	  \tt {\tiny [A$\langle$GM]}  
	 Your message has been relayed to the other player.\\ \tt \\ \tt \\ \tt You are currently playing round 9 of maximum 20.\\ \tt You have collectively accumulated 3 of 12 penalties.\\ \tt The other player sent this message:\\ \tt \\ \tt \\ \tt "(7,7) is empty for me. I will move H to (7,7). Confirm?"\\ \tt \\ \tt \\ \tt What is your next command? 
	  } 
	   } 
	   } 
	 & \\ 
 

    \theutterance \stepcounter{utterance}  

    & \multicolumn{4}{p{0.6\linewidth}}{\cellcolor[rgb]{0.8,1,0.9}{%
	\makecell[{{p{\linewidth}}}]{% 
	  \tt {\tiny [A$\rangle$GM]}  
	 SAY: Confirmed. Moving H to (7,7) now. 
	  } 
	   } 
	   } 
	 & & \\ 
 

    \theutterance \stepcounter{utterance}  

    & & & \multicolumn{2}{p{0.3\linewidth}}{\cellcolor[rgb]{0.95,0.95,0.95}{%
	\makecell[{{p{\linewidth}}}]{% 
	  \tt {\tiny [GM$|$GM]}  
	 SAY: Confirmed. Moving H to (7,7) now. 
	  } 
	   } 
	   } 
	 & & \\ 
 

    \theutterance \stepcounter{utterance}  

    & & \multicolumn{4}{p{0.6\linewidth}}{\cellcolor[rgb]{0.9,0.9,0.9}{%
	\makecell[{{p{\linewidth}}}]{% 
	  \tt {\tiny [GM$\rangle$B]}  
	 Your message has been relayed to the other player.\\ \tt \\ \tt \\ \tt You are currently playing round 10 of maximum 20.\\ \tt You have collectively accumulated 3 of 12 penalties.\\ \tt The other player sent this message:\\ \tt \\ \tt \\ \tt "Confirmed. Moving H to (7,7) now."\\ \tt \\ \tt \\ \tt What is your next command? 
	  } 
	   } 
	   } 
	 & \\ 
 

    \theutterance \stepcounter{utterance}  

    & & & \multicolumn{4}{p{0.6\linewidth}}{\cellcolor[rgb]{1,0.85,0.72}{%
	\makecell[{{p{\linewidth}}}]{% 
	  \tt {\tiny [GM$\langle$B]}  
	 MOVE: H, (7,7) 
	  } 
	   } 
	   } 
	  \\ 
 

    \theutterance \stepcounter{utterance}  

    & & & \multicolumn{2}{p{0.3\linewidth}}{\cellcolor[rgb]{0.95,0.95,0.95}{%
	\makecell[{{p{\linewidth}}}]{% 
	  \tt {\tiny [GM$|$GM]}  
	 MOVE: H, (7,7) 
	  } 
	   } 
	   } 
	 & & \\ 
 

    \theutterance \stepcounter{utterance}  

    & & & \multicolumn{2}{p{0.3\linewidth}}{\cellcolor[rgb]{0.95,0.95,0.95}{%
	\makecell[{{p{\linewidth}}}]{% 
	  \tt {\tiny [GM$|$GM]}  
	 Moved 'H' to (7,7) successfully. Your updated grid looks like this:\\ \tt ```\\ \tt  1234567\\ \tt ╔═══╤═╤═╗\\ \tt ║◌◌◌│C│◌║ 1\\ \tt ║◌W◌│◌├─╢ 2\\ \tt ║◌◌I├─┘◌║ 3\\ \tt ║◌◌◌│◌◌◌║ 4\\ \tt ╟─┬─┼───╢ 5\\ \tt ║◌│◌│◌T◌║ 6\\ \tt ║◌│◌│◌◌H║ 7\\ \tt ╚═╧═╧═══╝\\ \tt ```\\ \tt  
	  } 
	   } 
	   } 
	 & & \\ 
 

    \theutterance \stepcounter{utterance}  

    & & \multicolumn{4}{p{0.6\linewidth}}{\cellcolor[rgb]{0.9,0.9,0.9}{%
	\makecell[{{p{\linewidth}}}]{% 
	  \tt {\tiny [A$\langle$GM]}  
	 Your message has been relayed to the other player.\\ \tt \\ \tt \\ \tt You are currently playing round 10 of maximum 20.\\ \tt You have collectively accumulated 3 of 12 penalties.\\ \tt The other player moved an object on their grid.\\ \tt \\ \tt \\ \tt What is your next command? 
	  } 
	   } 
	   } 
	 & \\ 
 

    \theutterance \stepcounter{utterance}  

    & \multicolumn{4}{p{0.6\linewidth}}{\cellcolor[rgb]{0.8,1,0.9}{%
	\makecell[{{p{\linewidth}}}]{% 
	  \tt {\tiny [A$\rangle$GM]}  
	 MOVE: H, (7, 7) 
	  } 
	   } 
	   } 
	 & & \\ 
 

    \theutterance \stepcounter{utterance}  

    & & & \multicolumn{2}{p{0.3\linewidth}}{\cellcolor[rgb]{0.95,0.95,0.95}{%
	\makecell[{{p{\linewidth}}}]{% 
	  \tt {\tiny [GM$|$GM]}  
	 MOVE: H, (7, 7) 
	  } 
	   } 
	   } 
	 & & \\ 
 

    \theutterance \stepcounter{utterance}  

    & & & \multicolumn{2}{p{0.3\linewidth}}{\cellcolor[rgb]{0.95,0.95,0.95}{%
	\makecell[{{p{\linewidth}}}]{% 
	  \tt {\tiny [GM$|$GM]}  
	 Penalty: (7,7) is not empty, but contains 'H'.\\ \tt \\ \tt You have collectively accumulated 4 of 12 penalties.\\ \tt Please try again! 
	  } 
	   } 
	   } 
	 & & \\ 
 

    \theutterance \stepcounter{utterance}  

    & & \multicolumn{4}{p{0.6\linewidth}}{\cellcolor[rgb]{0.9,0.9,0.9}{%
	\makecell[{{p{\linewidth}}}]{% 
	  \tt {\tiny [A$\langle$GM]}  
	 Penalty: (7,7) is not empty, but contains 'H'.\\ \tt \\ \tt You have collectively accumulated 4 of 12 penalties.\\ \tt Please try again! 
	  } 
	   } 
	   } 
	 & \\ 
 

    \theutterance \stepcounter{utterance}  

    & \multicolumn{4}{p{0.6\linewidth}}{\cellcolor[rgb]{0.8,1,0.9}{%
	\makecell[{{p{\linewidth}}}]{% 
	  \tt {\tiny [A$\rangle$GM]}  
	 SAY: H is already at (7,7) on my grid. So all my objects are now at target positions: C(1,1), W(2,2), I(3,3), T(6,6), H(7,7). Are you finished? 
	  } 
	   } 
	   } 
	 & & \\ 
 

    \theutterance \stepcounter{utterance}  

    & & & \multicolumn{2}{p{0.3\linewidth}}{\cellcolor[rgb]{0.95,0.95,0.95}{%
	\makecell[{{p{\linewidth}}}]{% 
	  \tt {\tiny [GM$|$GM]}  
	 SAY: H is already at (7,7) on my grid. So all my objects are now at target positions: C(1,1), W(2,2), I(3,3), T(6,6), H(7,7). Are you finished? 
	  } 
	   } 
	   } 
	 & & \\ 
 

    \theutterance \stepcounter{utterance}  

    & & \multicolumn{4}{p{0.6\linewidth}}{\cellcolor[rgb]{0.9,0.9,0.9}{%
	\makecell[{{p{\linewidth}}}]{% 
	  \tt {\tiny [GM$\rangle$B]}  
	 Moved 'H' to (7,7) successfully. Your updated grid looks like this:\\ \tt ```\\ \tt  1234567\\ \tt ╔═══╤═╤═╗\\ \tt ║◌◌◌│C│◌║ 1\\ \tt ║◌W◌│◌├─╢ 2\\ \tt ║◌◌I├─┘◌║ 3\\ \tt ║◌◌◌│◌◌◌║ 4\\ \tt ╟─┬─┼───╢ 5\\ \tt ║◌│◌│◌T◌║ 6\\ \tt ║◌│◌│◌◌H║ 7\\ \tt ╚═╧═╧═══╝\\ \tt ```\\ \tt You are currently playing round 11 of maximum 20.\\ \tt You have collectively accumulated 4 of 12 penalties.\\ \tt The other player sent this message:\\ \tt \\ \tt \\ \tt "H is already at (7,7) on my grid. So all my objects are now at target positions: C(1,1), W(2,2), I(3,3), T(6,6), H(7,7). Are you finished?"\\ \tt \\ \tt \\ \tt What is your next command? 
	  } 
	   } 
	   } 
	 & \\ 
 

    \theutterance \stepcounter{utterance}  

    & & & \multicolumn{4}{p{0.6\linewidth}}{\cellcolor[rgb]{1,0.85,0.72}{%
	\makecell[{{p{\linewidth}}}]{% 
	  \tt {\tiny [GM$\langle$B]}  
	 SAY: Almost! My C is at (2,1) but should be at (1,1). I'll move C to (1,1) now. Then I'll be finished. 
	  } 
	   } 
	   } 
	  \\ 
 

    \theutterance \stepcounter{utterance}  

    & & & \multicolumn{2}{p{0.3\linewidth}}{\cellcolor[rgb]{0.95,0.95,0.95}{%
	\makecell[{{p{\linewidth}}}]{% 
	  \tt {\tiny [GM$|$GM]}  
	 SAY: Almost! My C is at (2,1) but should be at (1,1). I'll move C to (1,1) now. Then I'll be finished. 
	  } 
	   } 
	   } 
	 & & \\ 
 

    \theutterance \stepcounter{utterance}  

    & & \multicolumn{4}{p{0.6\linewidth}}{\cellcolor[rgb]{0.9,0.9,0.9}{%
	\makecell[{{p{\linewidth}}}]{% 
	  \tt {\tiny [A$\langle$GM]}  
	 Your message has been relayed to the other player.\\ \tt \\ \tt \\ \tt You are currently playing round 11 of maximum 20.\\ \tt You have collectively accumulated 4 of 12 penalties.\\ \tt The other player sent this message:\\ \tt \\ \tt \\ \tt "Almost! My C is at (2,1) but should be at (1,1). I'll move C to (1,1) now. Then I'll be finished."\\ \tt \\ \tt \\ \tt What is your next command? 
	  } 
	   } 
	   } 
	 & \\ 
 

    \theutterance \stepcounter{utterance}  

    & \multicolumn{4}{p{0.6\linewidth}}{\cellcolor[rgb]{0.8,1,0.9}{%
	\makecell[{{p{\linewidth}}}]{% 
	  \tt {\tiny [A$\rangle$GM]}  
	 SAY: Finished! All my objects are aligned. Let's end the game. 
	  } 
	   } 
	   } 
	 & & \\ 
 

    \theutterance \stepcounter{utterance}  

    & & & \multicolumn{2}{p{0.3\linewidth}}{\cellcolor[rgb]{0.95,0.95,0.95}{%
	\makecell[{{p{\linewidth}}}]{% 
	  \tt {\tiny [GM$|$GM]}  
	 SAY: Finished! All my objects are aligned. Let's end the game. 
	  } 
	   } 
	   } 
	 & & \\ 
 

    \theutterance \stepcounter{utterance}  

    & & \multicolumn{4}{p{0.6\linewidth}}{\cellcolor[rgb]{0.9,0.9,0.9}{%
	\makecell[{{p{\linewidth}}}]{% 
	  \tt {\tiny [GM$\rangle$B]}  
	 Your message has been relayed to the other player.\\ \tt \\ \tt \\ \tt You are currently playing round 12 of maximum 20.\\ \tt You have collectively accumulated 4 of 12 penalties.\\ \tt The other player sent this message:\\ \tt \\ \tt \\ \tt "Finished! All my objects are aligned. Let's end the game."\\ \tt \\ \tt \\ \tt What is your next command? 
	  } 
	   } 
	   } 
	 & \\ 
 

    \theutterance \stepcounter{utterance}  

    & & & \multicolumn{4}{p{0.6\linewidth}}{\cellcolor[rgb]{1,0.85,0.72}{%
	\makecell[{{p{\linewidth}}}]{% 
	  \tt {\tiny [GM$\langle$B]}  
	 MOVE: C, (1,1) 
	  } 
	   } 
	   } 
	  \\ 
 

    \theutterance \stepcounter{utterance}  

    & & & \multicolumn{2}{p{0.3\linewidth}}{\cellcolor[rgb]{0.95,0.95,0.95}{%
	\makecell[{{p{\linewidth}}}]{% 
	  \tt {\tiny [GM$|$GM]}  
	 MOVE: C, (1,1) 
	  } 
	   } 
	   } 
	 & & \\ 
 

    \theutterance \stepcounter{utterance}  

    & & & \multicolumn{2}{p{0.3\linewidth}}{\cellcolor[rgb]{0.95,0.95,0.95}{%
	\makecell[{{p{\linewidth}}}]{% 
	  \tt {\tiny [GM$|$GM]}  
	 Moved 'C' to (1,1) successfully. Your updated grid looks like this:\\ \tt ```\\ \tt  1234567\\ \tt ╔═══╤═╤═╗\\ \tt ║C◌◌│◌│◌║ 1\\ \tt ║◌W◌│◌├─╢ 2\\ \tt ║◌◌I├─┘◌║ 3\\ \tt ║◌◌◌│◌◌◌║ 4\\ \tt ╟─┬─┼───╢ 5\\ \tt ║◌│◌│◌T◌║ 6\\ \tt ║◌│◌│◌◌H║ 7\\ \tt ╚═╧═╧═══╝\\ \tt ```\\ \tt  
	  } 
	   } 
	   } 
	 & & \\ 
 

    \theutterance \stepcounter{utterance}  

    & & \multicolumn{4}{p{0.6\linewidth}}{\cellcolor[rgb]{0.9,0.9,0.9}{%
	\makecell[{{p{\linewidth}}}]{% 
	  \tt {\tiny [A$\langle$GM]}  
	 Your message has been relayed to the other player.\\ \tt \\ \tt \\ \tt You are currently playing round 12 of maximum 20.\\ \tt You have collectively accumulated 4 of 12 penalties.\\ \tt The other player moved an object on their grid.\\ \tt \\ \tt \\ \tt What is your next command? 
	  } 
	   } 
	   } 
	 & \\ 
 

    \theutterance \stepcounter{utterance}  

    & \multicolumn{4}{p{0.6\linewidth}}{\cellcolor[rgb]{0.8,1,0.9}{%
	\makecell[{{p{\linewidth}}}]{% 
	  \tt {\tiny [A$\rangle$GM]}  
	 SAY: finished? 
	  } 
	   } 
	   } 
	 & & \\ 
 

    \theutterance \stepcounter{utterance}  

    & & & \multicolumn{2}{p{0.3\linewidth}}{\cellcolor[rgb]{0.95,0.95,0.95}{%
	\makecell[{{p{\linewidth}}}]{% 
	  \tt {\tiny [GM$|$GM]}  
	 SAY: finished? 
	  } 
	   } 
	   } 
	 & & \\ 
 

    \theutterance \stepcounter{utterance}  

    & & \multicolumn{4}{p{0.6\linewidth}}{\cellcolor[rgb]{0.9,0.9,0.9}{%
	\makecell[{{p{\linewidth}}}]{% 
	  \tt {\tiny [GM$\rangle$B]}  
	 Moved 'C' to (1,1) successfully. Your updated grid looks like this:\\ \tt ```\\ \tt  1234567\\ \tt ╔═══╤═╤═╗\\ \tt ║C◌◌│◌│◌║ 1\\ \tt ║◌W◌│◌├─╢ 2\\ \tt ║◌◌I├─┘◌║ 3\\ \tt ║◌◌◌│◌◌◌║ 4\\ \tt ╟─┬─┼───╢ 5\\ \tt ║◌│◌│◌T◌║ 6\\ \tt ║◌│◌│◌◌H║ 7\\ \tt ╚═╧═╧═══╝\\ \tt ```\\ \tt You are currently playing round 13 of maximum 20.\\ \tt You have collectively accumulated 4 of 12 penalties.\\ \tt The other player sent this message:\\ \tt \\ \tt \\ \tt "finished?"\\ \tt \\ \tt \\ \tt What is your next command? 
	  } 
	   } 
	   } 
	 & \\ 
 

    \theutterance \stepcounter{utterance}  

    & & & \multicolumn{4}{p{0.6\linewidth}}{\cellcolor[rgb]{1,0.85,0.72}{%
	\makecell[{{p{\linewidth}}}]{% 
	  \tt {\tiny [GM$\langle$B]}  
	 SAY: finished! 
	  } 
	   } 
	   } 
	  \\ 
 

    \theutterance \stepcounter{utterance}  

    & & & \multicolumn{2}{p{0.3\linewidth}}{\cellcolor[rgb]{0.95,0.95,0.95}{%
	\makecell[{{p{\linewidth}}}]{% 
	  \tt {\tiny [GM$|$GM]}  
	 SAY: finished! 
	  } 
	   } 
	   } 
	 & & \\ 
 

    \theutterance \stepcounter{utterance}  

    & & & \multicolumn{2}{p{0.3\linewidth}}{\cellcolor[rgb]{0.95,0.95,0.95}{%
	\makecell[{{p{\linewidth}}}]{% 
	  \tt {\tiny [GM$|$GM]}  
	 true 
	  } 
	   } 
	   } 
	 & & \\ 
 

    \theutterance \stepcounter{utterance}  

    & & & \multicolumn{2}{p{0.3\linewidth}}{\cellcolor[rgb]{0.95,0.95,0.95}{%
	\makecell[{{p{\linewidth}}}]{% 
	  \tt {\tiny [GM$|$GM]}  
	 * success: True\\ \tt * lose: False\\ \tt * aborted: False\\ \tt -------\\ \tt * Shifts: 4.00\\ \tt * Max Shifts: 8.00\\ \tt * Min Shifts: 4.00\\ \tt * End Distance Sum: 0.00\\ \tt * Init Distance Sum: 23.66\\ \tt * Expected Distance Sum: 20.95\\ \tt * Penalties: 4.00\\ \tt * Max Penalties: 12.00\\ \tt * Rounds: 13.00\\ \tt * Max Rounds: 20.00\\ \tt * Object Count: 5.00\\ \tt  
	  } 
	   } 
	   } 
	 & & \\ 
 

    \theutterance \stepcounter{utterance}  

    & & & \multicolumn{2}{p{0.3\linewidth}}{\cellcolor[rgb]{0.95,0.95,0.95}{%
	\makecell[{{p{\linewidth}}}]{% 
	  \tt {\tiny [GM$|$GM]}  
	 * Main Score: 64.59\\ \tt \\ \tt -------\\ \tt * Distance Score: 1.00\\ \tt * Consistency Score: 1.00\\ \tt * Coverage Score: 0.56\\ \tt * Penalty Score: 0.82\\ \tt * Alternative Penalty Score: 0.90\\ \tt * Alternative Main Score: 90.00\\ \tt \\ \tt -------\\ \tt * Shifts: 4.00\\ \tt * Max Shifts: 8.00\\ \tt * Min Shifts: 4.00\\ \tt * End Distance Sum: 0.00\\ \tt * Init Distance Sum: 23.66\\ \tt * Expected Distance Sum: 20.95\\ \tt * Penalties: 4.00\\ \tt * Max Penalties: 12.00\\ \tt * Rounds: 13.00\\ \tt * Max Rounds: 20.00\\ \tt * Object Count: 5.00\\ \tt  
	  } 
	   } 
	   } 
	 & & \\ 
 

\end{supertabular}
}

\end{document}
